\documentclass[11pt,a4paper]{article}
\usepackage[T1]{fontenc}
\usepackage[utf8]{inputenc}
\usepackage{lmodern,amsmath,amssymb,amsthm,mathtools}
\usepackage[margin=25mm]{geometry}
\usepackage{booktabs,longtable,array,enumitem,microtype}
\usepackage[hidelinks]{hyperref}
\setlength{\parindent}{0pt}
\setlength{\parskip}{5pt plus 1pt}
\setlist{leftmargin=1.7em,itemsep=3pt,topsep=4pt}
\newtheorem{theorem}{Theorem}[section]
\newtheorem{proposition}[theorem]{Proposition}
\newtheorem{lemma}[theorem]{Lemma}
\newtheorem{corollary}[theorem]{Corollary}
\theoremstyle{definition}
\newtheorem{definition}[theorem]{Definition}
\newtheorem{example}[theorem]{Example}
\newtheorem{remark}[theorem]{Remark}
\newcommand{\Set}{\mathbf{Set}}
\newcommand{\Cat}{\mathbf{Cat}}
\newcommand{\FinBij}{\mathbf{FinBij}}
\newcommand{\V}{\mathcal V}
\newcommand{\Lcat}{\mathcal L}
\newcommand{\K}{\mathcal K}
\newcommand{\J}{\mathcal J}
\newcommand{\C}{\mathcal C}
\newcommand{\W}{\mathcal W}
\newcommand{\Prob}{\mathcal P}
\newcommand{\id}{\mathrm{id}}
\newcommand{\st}{\operatorname{St}}
\newcommand{\List}{\operatorname{List}}
\newcommand{\supp}{\operatorname{supp}}
\newcommand{\inl}{\mathrm{inl}}
\newcommand{\inr}{\mathrm{inr}}
\newcommand{\diamondc}{\mathbin{\diamond}}
\newcommand{\Sub}{\operatorname{Sub}}
\newcommand{\Hom}{\operatorname{Hom}}
\title{External Framing and Representation Comparison\\[5pt]
\large An exact full Set action criterion and the retained receipt boundary}
\author{}
\date{1 November 2026}
\begin{document}
\maketitle
\begin{abstract}
The entire canonical O1 mixed construction admits a pure-compatible external
Set action with invertible parameter comparisons exactly when transport from
the empty content fibre to every content fibre is an equivalence.
Necessity follows from empty interfaces and the empty frame; sufficiency
coherently converts the independent content presentation to constant grades
and derives the mixed action from canonical Set-functor strengths.
The criterion does not require invertible extension unit or multiplication
bodies. Constant ordered polynomial content satisfies it, with every admitted
augmented rig context and every native cell retained. For the full
collection family, the probability specialization maps strongly to extensional
finite laws on ordered output lists, with an exact local connected-component
quotient and a positive-history-only support adjunction.
Structured equivalences of independent presentations
induce local mixed equivalences, with separate conditions for preservation
of additional framing structure.

The full retained receipt model has a different answer. Uniform receipt
framing is strong at each fixed frame and on the bijection core, but only lax
in general frame functions. A two-inclusion argument excludes every
full Set-parameter pseudoaction with the specified pure parameter components
and unit comparison, independently of the framing recipe. This obstruction
is invariant under pure-compatible biequivalence; an information-forgetting
quotient can avoid it precisely because it is not such an equivalence.
The report is an O2 acceptance candidate with an exact criterion on the
entire canonical O1 family and explicit positive instances, not a claim that
all O1 inputs satisfy that criterion or that the expanded programme is
complete. TeX sources are supplied without compilation.
\end{abstract}
\tableofcontents

\section{Scope and the exact action requirement}\label{sec:scope}

O2 of the expanded research contract~\cite{Contract} requests a declared
external monoidal interface structure and a compatible action on a chosen
nontrivial fragment. This report chooses $(\Set,\times,1)$ and proves an
exact existence criterion for its full pure-compatible action on the entire
canonical O1 mixed construction~\cite{O1}. It then verifies the criterion
and gives explicit formulas for constant ordered polynomial content, including
a substantive information-forgetting probability comparison on that full
collection family. The context input remains any coherent symmetric rig
admitted by O1, not just numerical weights or a terminal resource category.
The criterion classifies which O1 inputs admit the stated action; it does
not assert that every input passes.

Three different operations must remain distinct: the relative sequential
module tensor $\odot_{\Set}$; the internal resource tensor $\otimes$ of $\V$;
and the external interface product $X\times A$. The finite independent
occurrence operation built from $\otimes$ is also not the external action.
No binary tensor of effectful computations, interchange theorem, concurrent
scheduler, or commutation of two sequential effect orders is claimed.

Write $J(v)$ for a pure computation induced by a set map $v$ and write
$g\circ f$ for execution of $f$ followed by $g$. The strongest action
requirement used here comprises the following data.
\begin{enumerate}[label=\textup{(A\arabic*)}]
\item For every set $X$, functors
$\st_X:H(A,B)\to H(X\times A,X\times B)$ on the complete hom categories,
strong for the relative unit $j$ and multiplication $m$, and compatible with
all pure endpoint squares. In particular
$\st_XJ(v)\cong J(\id_X\times v)$ coherently for pure maps $v$.
\item Coherent unit and associativity comparisons
$\st_1\simeq\id$ and $\st_X\st_Y\simeq\st_{X\times Y}$, with the
canonical pure interface unitors and associators inserted.
\item For every function $u:X\to Y$, invertible comparison cells
\begin{equation}\label{eq:parameter}
 J(u\times\id_B)\circ\st_X(F)
 \;\cong\;\st_Y(F)\circ J(u\times\id_A),
\end{equation}
natural in all native $F$-cells, coherent under identity and composition
of $u$, and compatible with sequential composition and the comparisons
in the preceding items. Thus parameter naturality means horizontal
\emph{pseudonaturality}, not just existence of a native square.
\end{enumerate}
When the comparison in~\eqref{eq:parameter} is not required invertible we
will explicitly say \emph{lax parameter comparison}. Fixed-$X$ strong
functoriality alone does not imply (A3).

\begin{theorem}[Exact full Set action criterion]\label{thm:main}
For an entire canonical O1 mixed construction $H$, including the original
genuine one-sided relative pseudopromonad structure, conditions
\textup{(A1)--(A3)} can be supplied if and only if every canonical transport
\[
 E_B=(!_B)_!:\Lcat_\varnothing\longrightarrow\Lcat_B,
 \qquad !_B:\varnothing\longrightarrow B,
\]
is an equivalence. Equivalently, the intrinsic pure-postcomposition functors
$J(!_B)\circ-:H(\varnothing,\varnothing)\to H(\varnothing,B)$ are equivalences.
Empty frames and empty interfaces are included. No condition requires
$!_B\in\J$ or invertibility of the content extension bodies.
The criterion is invariant under base-object-identity, pure-compatible
biequivalence. Its sufficiency assertion is within the O1 class, not an
assertion about arbitrary bicategories.

In particular, fix any O1 coherent symmetric rig $\V$ and admitted map class $\J$ for
$T=\List$. Take the full constant category of ordered polynomial collections
$\C=\Set^{\mathbb N}$ as the content fibre, including all arity-preserving
native shape maps. The resulting retained mixed construction satisfies
\textup{(A1)--(A3)} for all sets and all frame functions. Its comparisons
preserve the O1 crossing and its four modifications. The action does not
change any context resource coefficient or content collection grade.

Structured equivalences of independent O1 inputs induce coherent local
mixed equivalences. On the positive frame fragment these are frame compatible
under the stated one-sided frame-preservation condition. Within the
nonnegative-semiring specialization of the full collection family, raw finite
histories admit a strong, frame-preserving map to finite-support laws on lists,
whose local discrete quotient is exactly connected-component quotienting.
This map is neither locally full nor locally faithful and does not reflect
invertibility.
\end{theorem}
The exact criterion is proved in Section~\ref{sec:exact}. The explicit
constant-polynomial action is proved directly in Section~\ref{sec:constant};
the representation assertions are proved in Sections~\ref{sec:equivalence}
and \ref{sec:probability}. Section~\ref{sec:transfer} first gives the
one-sided transfer mechanism used in the general sufficiency proof.
Sections~\ref{sec:receipts} and \ref{sec:nogo} isolate two different receipt
boundaries. The main theorem is a necessary-and-sufficient test, not an
unconditional full-Set action theorem for every O1 contentad.

\section{Inherited data and retained cells}\label{sec:data}

We use the universe convention and relative pseudocoend presentation of
O1~\cite{O1}. Interfaces and shape sets may be infinite; each context fibre
and each polynomial position set is finite. $\J$ contains isomorphisms and
is closed under composition and pullback. O1 additionally requires
$T(\J)\subseteq\J$. No projection $X\times B\to B$ is assumed in $\J$.

A context $P$ over $A$ has a finite carrier $P_a$ at each $a$, coefficients
$V_{a,p}\in\V$, and a global augmentation
$\epsilon_a:\bigoplus_{p\in P_a}V_{a,p}\to I$.
A native map has a carrier function $r_a:P_a\to P'_a$ and actual maps
\[
 \bigoplus_{r_a(p)=p'}V_{a,p}\longrightarrow V'_{a,p'}
\]
whose composite with the target augmentation is the source augmentation,
using the prescribed sum reassociation. These are not replaced by scalar
inequalities or carrier functions. Write
$e_AP=\coprod_{a\in A}P_a$ and $\pi_P:e_AP\to A$; admissibility means
$\pi_P\in\J$. The rig sums are monoidal sums and need not be categorical
coproducts. O1 supplies dependent accumulation, occurrence tensors, their
actual aggregate coefficient maps, and their coherence.

A general contentad has a total category $p:\Lcat\to\Set$, extension
$d:\Lcat\to\Set$, unit grades $U_B$, accumulation $Q\diamondc R$ for
$pR=dQ$, and bodies
\[
 \eta^L_B:B\to dU_B,\qquad
 \mu^L_{Q,R}:dR\to d(Q\diamondc R).
\]
We inherit its specified global opcartesian transports, genuine cartesian
lifts along $\J$, and occurrence map $\nu:d\Rightarrow Tp$.
The O1 display property says that, for $v:E\to B$ in $\J$,
\begin{equation}\label{eq:display}
 d(v^*Q)\cong dQ\mathbin{\times}_{TB}TE,
\end{equation}
with the actual cartesian projection and arity as its two coordinates.
The content-only unit, accumulation and Beck--Chevalley conditions, and
the full resulting mixed law, are dependencies, not reproved assertions.

The mixed hom category has objects
\[
 H(A,B)=(e_A\downarrow d_B),\qquad
 F=(P,Q,f:e_AP\to dQ).
\]
A native cell is $(\alpha:P\to P',\beta:Q\to Q')$ with
$d\beta\,f=f'e_A\alpha$.
For $v:B\to B'$ and $u:A\to A'$, an endpoint square is
\begin{equation}\label{eq:endpoint}
 \alpha:P\to u^*P',\quad \beta:v_!Q\to Q',\qquad
 d\beta\,d\kappa_{v,Q}f=f'e_u(P')e_A\alpha.
\end{equation}
Here $\kappa$ is the chosen opcartesian arrow. Equality of mixed cells is
equality of their actual context and grade components. Faithfulness of $d$
or of a coefficient-to-set functor is never presumed.

O1 constructs the crossing by pulling $Q$ back along the following
context projection $\pi:E\to B$. Over $y\in dQ$ it uses one independent
context occurrence for every addressed position of $\nu_Q(y)$, with the
original coefficients tensored over those positions. Display~\eqref{eq:display}
determines the crossing body from $y$ and the complete addressed witness.
O1 proves native-cell and endpoint naturality, pure-middle balance, four
invertible crossing modifications, and their ten coherence diagrams. Its
composite theorem supplies $j,m$ on $H$. We will check the new action on
those same components and modifications rather than infer it from an
unexamined monoidal symbol.

\section{Full Set framing for constant ordered polynomial content}\label{sec:constant}

\subsection{The chosen contentad and explicit sequential form}
For $C=(C_n)_{n\ge0}$ put
\[
 F_C(B)=\coprod_{n\ge0}C_n\times B^n.
\]
A native shape map $c:C\to C'$ is a family $c_n:C_n\to C'_n$.
The total category is $\Set\times\C$ and
$d(B,C)=F_C(B)$. A total arrow $(v,c)$ acts by
$(s;b_1,\ldots,b_n)\mapsto(c_n(s);v(b_1),\ldots,v(b_n))$.
Both endpoint reindexing and transport leave $C$ unchanged. The occurrence
map forgets $s$ and retains the ordered list. The unit collection $\mathbb I$
has a single unary shape and no other shapes.

Let $\Sub(C,D)$ be the collection representing $F_C(F_D(-))$.
Its shapes are a shape $s\in C_n$ and a tuple
$(t_i\in D_{m_i})_{1\le i\le n}$; its positions are the ordered addresses
$(i,j)$ with $1\le j\le m_i$. This notation avoids reversing the
contentad's convention: if $C\diamondc D$ represents $F_DF_C$, then
$\Sub(C,D)=D\diamondc C$. Native maps act on $s$ and every $t_i$.
The substitution body is the evident addressed-list bijection. Nested
substitution is flattened by concatenating addresses, and units delete
unary unit nodes. These operations give the O1 constant contentad. Display
holds because the fixed shape and its ordered lifted list uniquely specify
a lifted extension point; transport comparisons and Beck--Chevalley are
identity on collection grades. Thus all O1 hypotheses are independently
verified for this content, not inferred from the desired action.

For clarity the inherited mixed composition can be written without relative
tensor notation. Let $F=(P,C,f):A\to B$ and $G=(R,D,g):B\to E$. Write
\[
 f(a,p)=(s;b_1,\ldots,b_n),\qquad
 g(b_i,r_i)=(t_i;e_{i1},\ldots,e_{im_i}).
\]
A composite context witness at $a$ is
\begin{equation}\label{eq:composite-witness}
 (p,(r_i)_{i=1}^n),\qquad r_i\in R_{b_i},
\end{equation}
with coefficient
\begin{equation}\label{eq:composite-coeff}
 V_{a,p}\otimes\bigotimes_{i=1}^n W_{b_i,r_i}.
\end{equation}
Its grade is $\Sub(C,D)$ and its body is
\begin{equation}\label{eq:composite-body}
 ((s,(t_i)_i);(e_{ij})_{(i,j)}).
\end{equation}
Empty position families have one empty witness and coefficient $I$.
Augmentation and all native maps are the O1 dependent accumulation and
occurrence aggregate maps. These formulas follow by its crossing, then
context accumulation, then content substitution. In particular every shape
of $\Sub(C,D)$ is retained, including shapes not reached by this body.

\subsection{Definition on objects and every native square}
For any set $X$, put $P^X_{(x,a)}=P_a$, with exactly the same coefficient
objects, augmentations and fibre maps. Its extension is canonically
$X\times e_AP$, and its carrier projection is $\id_X\times\pi_P$.
This is a pullback of $\pi_P$, so is in $\J$, including when $X$ is empty
or infinite. There is no need to pull a grade along $X\times B\to B$.

Define the canonical body strength
\begin{equation}\label{eq:poly-strength}
 b_{X,C,B}(x,(s;b_1,\ldots,b_n))
       =(s;(x,b_1),\ldots,(x,b_n)),
\end{equation}
and set
\[
 \st_X(P,C,f)=(P^X,C,b_{X,C,B}(\id_X\times f)).
\]
The grade is $C$, not $X\times C$. A native cell $(\alpha,c)$ is sent to
$(\alpha^X,c)$, where every coefficient map of $\alpha^X$ is the original
one. Naturality of~\eqref{eq:poly-strength} in all $c_n$ and endpoint maps
makes the body equation hold. Identity and composite cells are preserved
because their carrier maps, coefficient maps and shape maps are preserved.
This includes noninjective shape merges and noninvertible coefficient maps.

For an endpoint square~\eqref{eq:endpoint}, the framed square is over
$(\id_X\times u,\id_X\times v)$, has the same shape component and the
repeated context component, and satisfies its body equation because
\[
 F_c(\id_X\times v)b_{X,C,B}
   =b_{X,C',B'}(\id_X\times F_c(v)).
\]
Context reindexing commutes with framing fibrewise. Hence this is natural
on all endpoint squares, without assuming $u$ or $v$ in $\J$.

\subsection{Derivation of the crossing comparison}
Let $R$ be a context over $B$ with projection $\pi:E_R\to B$.
The crossed grade is still $C$. The occurrence context over
$y=(s;b_1,\ldots,b_n)\in F_CB$ has carrier
$\prod_iR_{b_i}$ and coefficient $\bigotimes_iW_{b_i,r_i}$.
After framing, the following two constructions agree:
\begin{enumerate}[label=\textup{(\roman*)}]
\item form this occurrence context and then repeat it at each $x$;
\item form the occurrence context of $R^X$ over $F_C(X\times B)$,
then reindex it along $b_{X,C,B}:X\times F_CB\to F_C(X\times B)$.
\end{enumerate}
At $(x,y)$ both carriers are $\prod_iR_{b_i}$ and both coefficients are
$\bigotimes_iW_{b_i,r_i}$. Their augmentations are the same occurrence
augmentation. The crossed bodies on either route send the tuple to
$(s;((x,b_i),r_i)_i)$ in $F_C(X\times E_R)$, under the evident identification
of the extension of $R^X$ with $X\times E_R$.
The grade comparison is the identity on $C$, and this body equality is also
an instance of display uniqueness.

This comparison is deliberately only over $X\times F_CB$, or after
reindexing along the actual framed body. It is not an isomorphism of the
unrestricted occurrence contexts over all $F_C(X\times B)$:
multicoloured lists need not lie in the image of $b$. At a nullary shape
$b$ also need not remember $x$. Neither surjectivity nor injectivity of $b$
is needed.

To verify naturality on native context maps, fix an occurrence tuple in the
target. The actual source aggregate map has the form
\[
 \bigoplus_{(r_i):\,\phi_i(r_i)=r'_i}\bigotimes_iW_{b_i,r_i}
 \longrightarrow
 \bigotimes_i\Bigl(\bigoplus_{\phi_i(r_i)=r'_i}W_{b_i,r_i}\Bigr)
 \longrightarrow\bigotimes_iW'_{b_i,r'_i},
\]
with the prescribed distributivity and tensor of native maps.
Both framed routes have this very same indexing set, the same primitive
maps, and the same structural maps in $\V$. Equality therefore holds in
$\V$, not merely on carriers. Shape maps preserve the ordered number of
positions, so they give exactly the same assertion on each shape. Pure
endpoint transport and arbitrary frame maps only change the endpoint
labels, and commute with these tuple maps. This proves native and endpoint
naturality of the crossing comparison.

\subsection{Strong sequential comparisons and complete coherence}
The pure unit computation has singleton context coefficient $I$ and the
unary collection $\mathbb I$. Framing preserves it, with body
$(x,a)\mapsto(\mathbb I;(x,a))$. Thus the unit comparison is the canonical
singleton-context isomorphism and identity shape map.

For $F,G$ as in~\eqref{eq:composite-witness}, composing their framed
versions at input $(x,a)$ first produces $((x,b_i))_i$. Every subsequent
witness is therefore $r_i\in R_{b_i}$ at that same $x$. The multiplication
comparison
\begin{equation}\label{eq:strong-seq}
 \st_X(G)\circ\st_X(F)\xrightarrow{\sim}\st_X(G\circ F)
\end{equation}
forgets only the redundant repeated frame coordinates in the notation.
It is a bijection on the complete witness tuples~\eqref{eq:composite-witness},
the identity on the collection $\Sub(C,D)$, and the identity on the
coefficient~\eqref{eq:composite-coeff}, up to the prescribed identical
bracketing. The body on both sides is~\eqref{eq:composite-body} with $x$
attached to every final leaf. This is an invertible native comparison,
including at all dormant shapes. Its naturality follows from the actual
aggregate-map check just made and naturality of shape substitution.

Here is a componentwise verification of all required coherence, not a
claim based solely on equal observations. For three composable computations,
write a full witness by its stage addresses
\[
 (p,(r_i)_i,(q_{ij})_{i,j}),
\]
retaining the corresponding nested shape tree and all primitive coefficient
objects. Both associativity paths for~\eqref{eq:strong-seq} erase repeated
copies of $x$ and retain this same addressed data. Their grade component is
the same tree-flattening bijection. Their coefficient component is the
inherited O1 associator/Fubini map on the same primitive objects in the same
address order; O1 coherence equates these actual maps. Their bodies attach
$x$ to the same final leaves. Thus the sequential associativity square
commutes. Either unit triangle deletes the same singleton witness and unary
unit node; its coefficient route is the inherited unit coherence, and its
body is unchanged.

The four crossing modifications are preserved individually. The context-unit
modification inserts or removes singleton occurrence witnesses; framing leaves
those witnesses and their $I$ coefficients unchanged. The content-unit
modification restricts to the unary unit shape; framing has exactly one
unchanged context occurrence. The context-accumulation modification is
dependent Fubini on the addressed choices; all primitive resource maps and
position shuffles are unchanged. The content-accumulation modification
flattens the same ordered shape tree; the frame on every leaf is $x$ before
and after flattening, and its grade map is the original flattening map.
Hence each modification square commutes on context, grade and body.

Every crossing coherence diagram is a pasting of those four modifications,
O1 structural maps and native naturality. After the above comparisons its
context component is the corresponding O1 diagram on unchanged coefficients;
its grade component is the corresponding tree reassociation/unit diagram;
its body is the same leaf-labelled map. The inherited ten O1 diagrams
therefore prove their framed versions. Likewise pure-middle balance has
unchanged native components and its middle pure map becomes
$\id_X\times v$; functoriality of endpoint labels proves its defining
relations. This verifies compatibility on the relative pseudocoend
presentation, rather than only on representatives with a fixed middle set.

\subsection{Frame parameters and action constraints}
For $u:X\to Y$, compare the two sides of~\eqref{eq:parameter}.
At input $(x,a)$ both have precisely the witnesses $P_a$, the same
coefficients and collection $C$, and the body
\[
 (s;(u(x),b_1),\ldots,(u(x),b_n)).
\]
The comparison is identity on these native components after deletion of
pure singleton units. It is invertible even when $u$ is neither injective
nor surjective. Naturality on cells follows because its context and grade
components are identities. For $v:Y\to Z$, the successive comparisons
and the comparison for $vu$ both replace the frame label by $v(u(x))$;
all their native components are the same unit identifications. This proves
parameter identity and composition coherence. Compatibility with
\eqref{eq:strong-seq} is checked on~\eqref{eq:composite-witness}:
both paths keep its witnesses and coefficients and put $u(x)$ on all
final leaves. The grade map remains identity on $\Sub(C,D)$.

The action unitor uses $1\times A\cong A$ and the corresponding body
identification; its native context and grade components are identities.
The action associator identifies $X\times(Y\times A)$ with
$(X\times Y)\times A$, retaining each witness and shape and replacing
$(x,(y,b))$ by $((x,y),b)$. Its naturality for arbitrary functions in either
parameter is equality of tuple functions. Its triangle and pentagon are
the ordinary product triangle and pentagon on labels, with identity
resource and shape maps. Compatibility of these constraints with the unit,
multiplication and crossing comparisons follows componentwise: both routes
retain the full addressed witness and shape tree, use the same O1 resource
map, and associate the frame tuple in the same prescribed way. This proves
all parts of Theorem~\ref{thm:main} concerning the action.

\begin{example}[A nonunit frame on an effectful computation]\label{ex:effect}
Let $X=\{\mathrm{red},\mathrm{blue}\}$ and let $C$ have one binary shape
$s$. A body returning $(s;b_0,b_1)$ is not the pure unary-unit body.
Its red frame returns $(s;(\mathrm{red},b_0),(\mathrm{red},b_1))$;
its blue frame returns the corresponding blue-labelled output. A following
context can choose genuinely different witnesses and coefficients at
$b_0,b_1$; the composition retains both addressed choices and their tensor
coefficients. Noninjective native maps merging two binary shapes, if present,
are still allowed. The action neither samples a fresh frame at each leaf nor
copies any resource coefficient. The frame is freely reusable classical
interface data, not a scarce capability object. A nullary shape remains a
valid effectful output but exposes no leaf label $x$; no recovery guarantee
is inferred.
\end{example}

\begin{corollary}[Fixed polynomial monad control]\label{cor:fixed}
For any O1 finitary cartesian polynomial monad $T$, the terminal-grade content
$d_B(*)=TB$ has a full Set action with
$f^X(x,z)=T(i_x)f(z)$, where $i_x(b)=(x,b)$.
\end{corollary}
\begin{proof}
The grade is unique. The occurrence witnesses and coefficients are those of
the polynomial shape of $f(z)$ and are unchanged. The crossing comparison
uniformly inserts $x$ at its addressed positions. Unit compatibility is
$T(i_x)\eta_B=\eta_{X\times B}i_x$; multiplication compatibility is
naturality of $\mu^T$ and the fact that successive occurrence substitutions
retain the same $x$. Every coefficient and addressed-position coherence is
exactly the one proved above, with the polynomial substitution position
bijection in place of ordered-list concatenation. Arbitrary parameter
functions have identity grade component and identical witnesses, so their
comparisons are invertible. Frame unit and associativity are functoriality
of $T$ on the corresponding insertions. This proof does not require
$\eta^T$ or $\mu^T$ invertible, and makes no claim that nullary outputs retain
an observable $x$.
\end{proof}

\section{A one sided criterion that derives mixed compatibility}\label{sec:transfer}

The direct proof above establishes the mandatory positive case without
assuming any mixed strength. The following extension explains exactly what
additional one-sided data can transfer the argument to another O1 contentad.
It is a sufficient criterion, not a consequence of O1 alone.

\begin{definition}[Independent content framing data]\label{def:frame-data}
In addition to the O1 data in Section~\ref{sec:data}, give:
\begin{enumerate}[label=\textup{(F\arabic*)}]
\item A functor $S:\Set\times\Lcat\to\Lcat$ with
$pS_XQ=X\times pQ$, with coherent action unit and associativity isomorphisms
$S_1\cong\id$, $S_XS_Y\cong S_{X\times Y}$ over the base product
constraints, natural in all total content arrows and frame maps.
\item Natural maps $b_{X,Q}:X\times dQ\to dS_XQ$ respecting these action
constraints. They need not be injective, surjective or invertible.
\item Each $S_X$ preserves the specified global opcartesian arrows and the
cartesian arrows over $\J$. Hence it has the canonical transport comparisons
for $\id_X\times v$, and these are coherent by universal uniqueness.
\item One-sided unit and accumulation isomorphisms
\[
 u_X:S_XU_B\cong U_{X\times B},\qquad
 a_X:S_XQ\diamondc(b_{X,Q})_!S_XR\cong S_X(Q\diamondc R)
 \quad(pR=dQ),
\]
with body equations
\begin{align}
 d(u_X)b_{X,U_B}(\id_X\times\eta^L_B)&=\eta^L_{X\times B},
                                                        \label{eq:frame-unit}\\
 d(a_X)\mu^L_{S_XQ,(b_{X,Q})_!S_XR}
 d\kappa_{b_{X,Q},S_XR}b_{X,R}
 &=b_{X,Q\diamondc R}(\id_X\times\mu^L_{Q,R}).
                                                        \label{eq:frame-accum}
\end{align}
These maps are natural in the total composable-pair category and frame
parameters. They obey the one-sided unit and associativity axioms: the two
successive applications of $a_X$ to a composable triple agree after the
original content associator and the canonical opcartesian comparisons;
applying $u_X$ to either unit input agrees with the original left or right
unit constraint. They also commute with the action unitor and associator
from (F1). All these equations use only $U,\diamondc$, their one-sided
constraints, $S,b,u,a$ and the transport comparisons, not contexts or a mixed
crossing.
\item The arity equation holds:
\begin{equation}\label{eq:frame-arity}
 \nu_{S_XQ}b_{X,Q}=\operatorname{st}^T_{X,pQ}
 (\id_X\times\nu_Q),\qquad
 \operatorname{st}^T(x,t)=T(i_x)t.
\end{equation}
\end{enumerate}
For full Set parameter pseudonaturality add the separate condition
\begin{enumerate}[label=\textup{(F6)},start=6]
\item For every $u:X\to Y$, the total arrow
$S(u,\id_Q):S_XQ\to S_YQ$ over $u\times\id_{pQ}$ is opcartesian.
Equivalently its canonical vertical factor
$(u\times\id_{pQ})_!S_XQ\to S_YQ$ is invertible.
\end{enumerate}
(F6) is not the preservation assertion in (F3): it concerns changing the
frame parameter, not transporting a content grade while $X$ is fixed.
\end{definition}

\begin{theorem}[One sided transfer]\label{thm:transfer}
Conditions \textup{(F1)--(F5)} induce strong fixed-frame functors on the full
O1 mixed construction, coherent action unit and associativity, and lax
parameter comparisons. They derive, rather than assume, compatibility with
the mixed crossing and its four modifications. With \textup{(F6)} all
parameter comparisons are invertible, giving \textup{(A1)--(A3)}. Without
\textup{(F6)}, bijective parameters still have invertible comparisons.
\end{theorem}
\begin{proof}
Define
\[
 \st_X(P,Q,f)=(P^X,S_XQ,b_{X,Q}(\id_X\times f)).
\]
On cells use the original coefficient maps repeated at $x$ and $S_X\beta$.
Naturality of $b$ proves the comma equation. On an endpoint square its grade
map is
\[
 (\id_X\times v)_!S_XQ\cong S_X(v_!Q)\xrightarrow{S_X\beta}S_XQ'.
\]
(F3) supplies the displayed isomorphism and its compatibility with the
opcartesian arrows; composing with those arrows and using $b$-naturality
proves~\eqref{eq:endpoint}. Uniqueness of opcartesian factorization proves
coherence under composed endpoint maps. Contexts commute with reindexing
fibrewise as before.

Let $c:\pi^*Q\to Q$ be the chosen cartesian arrow for a context projection
$\pi:E\to B$. Since $S_Xc$ is cartesian, there is a unique isomorphism
\begin{equation}\label{eq:cart-comparison}
 C_{X,\pi,Q}:S_X(\pi^*Q)\xrightarrow{\sim}
          (\id_X\times\pi)^*S_XQ
\end{equation}
whose composite with the target cartesian projection is $S_Xc$.
This is an actual grade isomorphism, not an isomorphism inferred from a
set-level bijection. Naturalities and iterated-lift coherence follow by
composing the candidate routes with that same projection and applying
cartesian uniqueness. Native context carrier maps are treated by the O1
cartesian/opcartesian factorization, with the comparisons from (F3).

Write $\widehat P$ for the original occurrence context over $dQ$ and
$\widehat{P^X}$ for the one over $dS_XQ$. Equation~\eqref{eq:frame-arity}
gives a native context isomorphism
\begin{equation}\label{eq:restricted-occurrence}
 b_{X,Q}^*\widehat{P^X}\cong\operatorname{pr}_{dQ}^*\widehat P
                 \quad\text{over }X\times dQ.
\end{equation}
At $(x,y)$ its addressed choices and primitive coefficients are identical:
the polynomial shape of $\nu_Q(y)$ has simply had $x$ attached at every
position. The native aggregation maps on both sides are the same finite
distributivity followed by the tensor of the same primitive maps. Thus
\eqref{eq:restricted-occurrence} is natural in all actual coefficient maps.
It is only claimed after this reindexing, not over all $dS_XQ$.

For the body, start with an original addressed witness $z$ above $y$.
Apply the original crossing body, then $b_{X,\pi^*Q}$ and $dC_{X,\pi,Q}$.
Its projection to $dS_XQ$ is $b_{X,Q}(x,y)$ by naturality of $b$ and the
defining triangle of~\eqref{eq:cart-comparison}. Its occurrence coordinate
in $T(X\times E)$ is the original witness with uniform frame $x$ by
\eqref{eq:frame-arity}. Those are precisely the two coordinates of the
framed crossing body. Display at $\id_X\times\pi$ identifies the two
bodies. This proves the derived crossing comparison and its naturality.

Context framing preserves its own unit and accumulation on identical
witnesses and coefficient maps. Equations~\eqref{eq:frame-unit} and
\eqref{eq:frame-accum} supply the strong content comparisons. Pasting
these with~\eqref{eq:cart-comparison} and
\eqref{eq:restricted-occurrence} gives an invertible mixed unit and
multiplication comparison. The latter is evaluated along the actual body
$b_{X,Q}(\id_X\times f)$, so no unrestricted occurrence comparison has
been used.

For completeness the four crossing-modification checks separate into native
components. For the context unit, both grade routes are the unique maps
between lifts of its unit projection, and both coefficient routes are the
same occurrence-unit map. For the content unit, $u_X$ identifies the grades,
\eqref{eq:frame-unit} identifies the bodies, and the arity-unit equation
identifies the same unary occurrence. For context accumulation, both grade
routes are cartesian lifts of the same composite carrier projection and
have the same composite with its original grade projection; cartesian
uniqueness identifies them. Their resource routes are the original dependent
Fubini maps on unchanged primitive coefficients. For content accumulation,
naturality of $a_X$ applied to the composable-pair lift makes the two grade
routes have the same projection to the framed accumulated grade. The
one-sided associativity and transport axioms in (F4), followed by cartesian
uniqueness, identify them. The arity accumulation equation identifies the
same flattened polynomial positions, so display identifies the bodies.

Every higher required pasting is now checked in three parts. Its context
component is the inherited O1 coherence on those primitive coefficients;
its grade component is (F1), (F3), (F4) coherence followed by cartesian or
opcartesian uniqueness; and its body component is the defining display
coordinate equality. This proves the action constraints, strong $j,m$
coherence and crossing coherence, including relative balance generators.
It does not cancel grade maps through a nonfaithful $d$.

For a frame map $u$, use the unchanged context coefficients and the vertical
factor of $S(u,\id_Q)$. Naturality of $b$ gives its square. Its composition,
sequential and action coherences follow respectively from functoriality of
$S$, the frame naturality in (F4), and (F1)--(F2). These are exactly the
lax parameter comparisons. Under (F6) their grade factors are invertible,
while their context components are the identity identifications, so the
comparisons are invertible. If $u$ is a bijection, $S(u,\id_Q)$ is itself an
isomorphism and therefore opcartesian. This proves all assertions.
\end{proof}

For constant polynomial content take $S_X(B,C)=(X\times B,C)$ and
$b$ as in~\eqref{eq:poly-strength}. The transport comparisons are identity
on $C$, (F4) is ordinary ordered shape substitution, and (F6) holds because
its grade map is identity. This independently verifies every hypothesis.
For terminal-grade fixed $T$, all grade comparisons are unique;
functoriality and naturality of $\eta^T,\mu^T$ give the body equations.
Canonical Set strength for any endofunctor $F$ is $(x,t)\mapsto F(i_x)t$,
and every natural transformation respects it. This fact supplies the body
identities but, by itself, does not supply O1 display or crossing hypotheses.

\subsection{A separate finite copower sufficient condition}
There is an independently checkable sufficient source of (F1)--(F5).
Assume the total content category has finite coproducts, $p,d$ preserve
them, the one-sided $U$ and accumulation functor preserve them, and $d$
reflects isomorphisms of vertical arrows. Coproducts in the composable-pair
category are componentwise using the preserved $p,d$ identifications.
For finite $X$ put $S_XQ=\coprod_{x\in X}Q$. Then
$pS_XQ\cong X\times pQ$ and $dS_XQ\cong X\times dQ$; use the latter
as $b$. Parameter maps are induced by the coproduct injections.

\begin{proposition}\label{prop:copower}
These conditions give strong finite fixed-frame functors and a coherent
strong $\FinBij$ action, with lax comparisons for arbitrary finite frame
functions. They do not in general imply \textup{(F6)}.
\end{proposition}
\begin{proof}
Finite copower universal properties give (F1)--(F2), including all action
constraints. A coproduct of opcartesian lifts is opcartesian over the
coproduct base map: a candidate map out of it factors uniquely on each
summand, and the factors assemble uniquely. This proves the opcartesian
part of (F3).

For a cartesian lift over $v:E\to B$, factor the coproduct of its arrows
through the chosen target cartesian lift to obtain a vertical map
\[
 c:\coprod_x v^*Q\longrightarrow
            (\id_X\times v)^*\Bigl(\coprod_xQ\Bigr).
\]
By display, its extension is the constant-frame insertion map
\[
 X\times(dQ\times_{TB}TE)
 \longrightarrow (X\times dQ)\times_{T(X\times B)}T(X\times E).
\]
It is bijective. An element on the right contains an explicit $x$; the
polynomial label forces every addressed frame to equal $x$. Delete those
frame entries for its inverse. At a nullary shape the explicit $x$ is still
in the first coordinate. Vertical conservativity of $d$ makes $c$ an
isomorphism, proving the cartesian part of (F3). Its naturality and iterated
coherence follow from its defining cartesian triangle.

Preservation of coproducts by $U$ and accumulation gives $u_X,a_X$.
Their body equations and one-sided coherences can be checked on each
coproduct injection, where they are the original unit, multiplication and
coherence equations. Joint universality of the injections proves the total
equalities. Naturality of $\nu$ on those injections gives (F5). Theorem
\ref{thm:transfer} applies. Nothing in this argument makes the fold or
inclusion induced by a nonbijective $X\to Y$ opcartesian, so it does not
establish (F6).
\end{proof}
This criterion is sufficient, not necessary. Constant polynomial evaluation
generally does not preserve these total coproducts, yet Section
\ref{sec:constant} already proves its full Set action.

\section{An exact criterion from the empty content fibre}\label{sec:exact}

We now prove the necessary-and-sufficient assertion in
Theorem~\ref{thm:main}. It concerns the entire canonical O1 construction,
with all original content grades and native arrows retained. The criterion
does not prescribe a frame recipe or assume a mixed strength. Its sufficient
direction uses the already valid one-sided \emph{pseudo} structure, including
invertible relative balance and associativity, not a merely lax operation
on grades.

\subsection{Necessity and the intrinsic formulation}
There is a canonical equivalence
\begin{equation}\label{eq:empty-hom}
 H(\varnothing,B)\simeq\Lcat_B.
\end{equation}
Indeed the context over the empty interface is unique, has empty extension,
and has only its identity context arrow. For each $Q\in\Lcat_B$ there is
one empty body, and its native arrows are exactly the grade arrows.
The empty context projection is $\id_\varnothing$, always admitted by
$\J$. Under~\eqref{eq:empty-hom}, pure postcomposition by $J(v)$ is the
original covariant transport $v_!$, through the original one-sided unit
comparison. In particular $J(!_B)\circ-$ is $E_B=(!_B)_!$.

Suppose (A1)--(A3) hold. Restrict $\st_\varnothing$ to these hom categories
and write $R_B:\Lcat_B\to\Lcat_\varnothing$ for the resulting functor.
Use parameter naturality for the unique frame map $\varnothing\to1$ on a
computation $\varnothing\to B$. Its source pure component is
$\id_\varnothing$, and its target pure component is $!_B$. The action
unitor and~\eqref{eq:parameter} therefore give a natural isomorphism
\begin{equation}\label{eq:empty-counit}
 E_BR_B\cong\id_{\Lcat_B}.
\end{equation}
For $B=\varnothing$, this gives $R_\varnothing\cong\id$, since
$E_\varnothing$ is canonically identity.

For $Q\in\Lcat_\varnothing$, realize $E_BQ$ by the composite
$J(!_B)\circ Q$. Strong sequential preservation and pure preservation give
\[
 R_BE_BQ\cong
 \st_\varnothing(J(!_B))\circ\st_\varnothing(Q)
 \cong J(\id_\varnothing)\circ R_\varnothing Q
 \cong Q.
\]
These are natural isomorphisms on all grade arrows. Thus $R_B$ is a
quasi-inverse to $E_B$. No formula for $\st_\varnothing$ on grades or
contexts was used. This proves necessity and the equivalence with the
intrinsic pure-postcomposition criterion.

\subsection{All map accumulation covariance from one sided validity}
The next lemma is needed for sufficiency. It is not an application of
cartesian preservation along $\J$, and it remains valid when
$!_B\notin\J$.

\begin{lemma}[Covariance of the existing accumulation]\label{lem:covariance}
For any set map $u:B\to C$, $Q\in\Lcat_B$ and $R\in\Lcat_{dQ}$, the
original one-sided relative pseudopromonad supplies canonical coherent
isomorphisms
\begin{equation}\label{eq:accum-covariance}
 u_!(Q\diamondc R)\cong
       (u_!Q)\diamondc(d\kappa_{u,Q})_!R,\qquad
 u_!U_B\cong U_C.
\end{equation}
They are natural in all native composable-pair arrows and compatible with
successive base transport, units and associative accumulation. Writing
$\gamma_{u,Q,R}$ for the first comparison, they satisfy the body equation
\begin{equation}\label{eq:covariance-body}
 d\gamma_{u,Q,R}\,d\kappa_{u,Q\diamondc R}\,\mu^L_{Q,R}
 =
 \mu^L_{u_!Q,(d\kappa_{u,Q})_!R}\,
 d\kappa_{d\kappa_{u,Q},R}.
\end{equation}
\end{lemma}
\begin{proof}
In the content-only module $N(A,B)=(A\downarrow d_B)$, take the arrow
$dR\to dQ$ represented by $R$ with identity body, followed by
$dQ\to B$ represented by $Q$ with identity body. Their composite has grade
$Q\diamondc R$. Postcompose with $J(u)$. Associating after the first
composition gives grade $u_!(Q\diamondc R)$. Associating $J(u)$ first
with the representative of $Q$ gives grade $u_!Q$ and body
$d\kappa_{u,Q}:dQ\to d(u_!Q)$; the remaining composition transports $R$
along that body and gives the other side of~\eqref{eq:accum-covariance}.
The original associator, pure transport and relative balance comparisons
are invertible. They produce a native isomorphism in $N(dR,C)$.
Its grade component is the required actual grade isomorphism, while its
comma body equation is exactly~\eqref{eq:covariance-body}. Thus both the
grade and extension-body comparisons are proved, without any faithfulness
assumption on $d$. Equivalently these are the output-endpoint
pseudonaturality comparisons of the original multiplication $m_N$ on
the two identity-body representatives.

Naturality is the native naturality of that associator and its balance
maps. For two maps $u,v$, the two transport comparisons are two pastings
of the original associator and pure-middle balance for the same string
of representatives and pure maps; the one-sided pentagon and balance
composition relation equate them. Adding another accumulation input gives
the same pentagon check with one additional content representative.
Inserting a unit gives the original triangle. Finally the relative unit
$j$ is covariant in its output endpoint; its comparison on a pure arrow
$u$ has grades $u_!U_B$ and $U_C$, giving the second isomorphism and its
unit/transport coherence. These arguments use only the independently valid
one-sided relative pseudopromonad.
\end{proof}

\subsection{Converting the independent content presentation to constant grades}
Assume every $E_B$ is an equivalence and put $\mathcal C_0=\Lcat_\varnothing$.
Normalize $E_\varnothing$ to the identity if necessary. Choose the
opcartesian arrows
$\kappa_{B,Q}:Q\to E_BQ$ over $!_B$.
There is a total functor over $\Set$
\begin{equation}\label{eq:constantization}
 E:\Set\times\mathcal C_0\longrightarrow\Lcat.
\end{equation}
For $f:B\to D$ and $a:Q\to Q'$ in $\mathcal C_0$, define $E(f,a)$ as the
unique arrow over $f$ satisfying
\[
 E(f,a)\kappa_{B,Q}=\kappa_{D,Q'}a.
\]
It exists and is unique by the opcartesianness of $\kappa_{B,Q}$, since
$f!_B=!_D$. This equation proves identities, composition and the native
naturality in $a$, so~\eqref{eq:constantization} is an actual functor.
Define honest Set endofunctors and native natural transformations by
\[
 F_Q(B)=d(E_BQ),\qquad F_Q(f)=dE(f,\id_Q),\qquad
 F_a(B)=dE_B(a).
\]
Thus $Q\mapsto F_Q$ is a functor
$\mathcal C_0\to[\Set,\Set]$. None of the maps $F_Q(f)$ is presumed
invertible.

The canonical opcartesian comparisons give $f_!E_B\cong E_D$. Consequently
every $f_!$ is an equivalence, since $E_B,E_D$ are equivalences. In an
opfibration with equivalent transition functors, an opcartesian arrow is
also cartesian. To check this assertion, let $c:Q\to f_!Q$ be opcartesian
over $f:B\to D$, and let an arrow into $f_!Q$ lie over $fg$ for
$g:A\to B$. Its vertical factor after the opcartesian lift over $fg$ is a
map $f_!g_!P\to f_!Q$, with the canonical transport identification.
Full faithfulness of $f_!$ uniquely lifts this to $g_!P\to Q$; its composite
with the lift over $g$ is exactly the required unique factor through $c$.
This proves the cartesian property against arbitrary candidate base arrows.

The functor $E$ preserves opcartesian arrows by construction. Cartesian
arrows in its constant source fibration have invertible grade component;
their images are compositions of an opcartesian arrow and a vertical
isomorphism and hence are cartesian by the preceding paragraph. In
particular it preserves the original $\J$-cartesian structure. Chosen
cartesian lifts are identified by their unique projection-preserving
isomorphisms, so no new choice of incompatible marks is made.

Fibrewise quasi-inverses extend coherently to an inverse over the fixed base:
for a total arrow over $f$, its arrow in the constant presentation is
obtained by opcartesian factorization and the full faithfulness of $E_D$.
Uniqueness makes these lifts functorial, and the fibre equivalence units and
counits give total units and counits. Thus $E$ is an equivalence of the
specified transport presentations, not only a separate equivalence of
unrelated fibres.

\subsection{The constant grade unit accumulation and bodies}
Define the constant unit to be $U_\varnothing$ and define the grade product
in the original contentad's $\diamondc$ convention:
\begin{equation}\label{eq:constant-product}
 Q\star R=Q\diamondc E_{F_Q(\varnothing)}R
                    \quad\text{in }\mathcal C_0.
\end{equation}
Apply Lemma~\ref{lem:covariance} to $!_B$ and this composable pair.
There is a canonical comparison
\[
 (d\kappa_{B,Q})_!E_{F_Q(\varnothing)}R
       \cong E_{F_Q(B)}R,
\]
because both are transport of $R$ along the unique map
$\varnothing\to F_Q(B)$:
$d\kappa_{B,Q}\,!_{F_Q(\varnothing)}=!_{F_Q(B)}$.
It follows that, naturally in $B,Q,R$,
\begin{equation}\label{eq:constant-accum}
 E_B(Q\star R)\cong
       E_BQ\diamondc E_{F_Q(B)}R,\qquad
 E_BU_\varnothing\cong U_B.
\end{equation}
The inherited extension bodies therefore become natural transformations
\begin{equation}\label{eq:constant-bodies}
 \eta:\id_{\Set}\longrightarrow F_{U_\varnothing},\qquad
 \mu_{Q,R}:F_R F_Q\longrightarrow F_{Q\star R}.
\end{equation}
For $\mu$, use the original accumulation body on the right side of
\eqref{eq:constant-accum}, then the inverse extension of that grade
isomorphism. Its naturality follows from native naturality of the original
body and the coherent covariance isomorphism. The same construction gives
the displayed $\eta$.

The product has coherent unit and associator isomorphisms inherited from
the one-sided contentad, as can be checked without inverting either body.
For the right unit, $E_{dQ}U_\varnothing\cong U_{dQ}$ makes
$Q\star U_\varnothing$ the original right-unit accumulation. For the left
unit, the transport of $Q$ along
$\eta_\varnothing:\varnothing\to dU_\varnothing$ is
$E_{dU_\varnothing}Q$, so the original left-unit comparison gives
$U_\varnothing\star Q\cong Q$.
For three grades, the original content associator compares the two nested
accumulations after transporting the third grade along the intervening
extension multiplication. But for every map $v:D\to D'$,
\[
 v_!E_D S\cong E_{D'}S.
\]
This applies even when $v$ is the noninvertible multiplication body in
\eqref{eq:constant-bodies}. Hence the two nested accumulations identify
with $(Q\star R)\star S$ and $Q\star(R\star S)$, giving the required
associator. Its pentagon, both unit triangles, and their body equations
are the original one-sided diagrams with the coherent comparisons of
Lemma~\ref{lem:covariance} inserted. The comparisons for any two iterated
base transports agree by opcartesian uniqueness, so these insertions do
not create a new coherence condition. This proves the constant one-sided
content structure, not just an objectwise product formula.

With extension $F_Q(B)$, the functor $E$ now has identity extension
comparison by definition and is a strong contentad equivalence through
\eqref{eq:constant-accum}. The original arity becomes the natural map
$\nu_{E_BQ}:F_Q(B)\to TB$. All arity equations, display squares,
$\J$-cartesian unit/accumulation preservation and Beck--Chevalley
comparisons transfer along the equivalence $E$, since it preserves the
specified cartesian and opcartesian arrows and their projections.
Thus the constantized presentation is itself an O1 input with the same
rig, map class, polynomial monad and retained native data up to the
specified structured equivalence.
Only the original display property along $\J$ is transferred here.
Existence of cartesian lifts along every base map does not assert display
exactness outside $\J$. The next construction uses Set functoriality,
not display or a cartesian lift along $X\times B\to B$.

\subsection{Canonical strength and sufficiency}
Leave the constant grade $Q$ unchanged under every frame, and define
\begin{equation}\label{eq:constantized-strength}
 b_{X,Q,B}(x,z)=F_Q(i_x)(z),\qquad i_x:B\to X\times B.
\end{equation}
Functoriality of $F_Q$ proves naturality in endpoint and frame maps and
the frame unit/associativity equations. Every natural transformation
$\alpha:F\to G$ respects this strength, since naturality at $i_x$ says
$\alpha_{X\times B}F(i_x)=G(i_x)\alpha_B$.
The strength on $F_RF_Q$ is the successive one: at fixed $x$,
\[
 F_R(b_{X,Q,B})F_R(i_x^{F_Q(B)})
   =F_R\bigl(b_{X,Q,B}i_x^{F_Q(B)}\bigr)
   =F_RF_Q(i_x).
\]
Consequently all native grade transformations, the unit and multiplication
in~\eqref{eq:constant-bodies}, and the arity transformation respect these
strengths. The one-sided accumulation comparisons are identity on
$Q\star R$, and their coherent body equations are these naturality
equations.

In the constant fibration, $S_X$ and every parameter arrow
$S(u,\id_Q)$ have unchanged grade components. They preserve cartesian and
opcartesian arrows, and $S(u,\id_Q)$ is opcartesian for every frame map.
This verifies every item (F1)--(F6), including its complete one-sided
unit/accumulation coherence, without guessing extra framing data.
Theorem~\ref{thm:transfer} derives the full mixed action on this constantized
O1 presentation.

The rig functor is identity, the polynomial-monad map is identity, and
$E$ with its constructed strong content structure and identity extension
comparison is a structured input equivalence. The O1 parameter theorem
and the local equivalence argument of Theorem~\ref{thm:equivalence} give
a pure-compatible mixed equivalence, with all native grades retained up
to equivalence. Transport the action back through a coherent inverse.
Its unit, associator, sequential and parameter comparisons remain
invertible; their axioms follow by applying the locally faithful
equivalence and the equivalence triangle identities to the known constant
presentation diagrams. The pure parameter components are identified with
$J(u\times\id)$ by pure compatibility. This proves sufficiency and completes
the exact assertion of Theorem~\ref{thm:main}.

\subsection{Invariance and boundaries of the criterion}
The condition can be stated entirely inside the mixed construction as
equivalence of $J(!_B)\circ-:H(\varnothing,\varnothing)\to
H(\varnothing,B)$. A base-object-identity, pure-compatible biequivalence
gives a square between this functor and its counterpart, commuting through
an invertible comparison; its two vertical hom functors are equivalences.
One horizontal functor is therefore an equivalence if and only if the other
is. This proves the asserted invariance of the criterion itself.
The sufficiency argument still requires an O1 presentation; the intrinsic
condition is not asserted sufficient for arbitrary bicategories.

Constant ordered polynomial and terminal-grade fixed-monad content satisfy
the criterion immediately because all $E_B$ are identity on their constant
grade categories. Receipts do not: $E_1$ transports labels from $T\varnothing$,
and the unary label $\eta_1(*)$ is not in the image of
$T(\varnothing\to1)$. Indeed the cartesian unit square at
$\varnothing\to1$ has empty upper-left set, so the fibre over
$\eta_1(*)$ is empty. A singleton receipt with that label is not isomorphic
to any transported grade. Thus essential surjectivity already fails.

Necessity uses the empty frame and empty interfaces. If either is excluded,
it is not an established converse for that changed problem. The independent
two-inclusion receipt obstruction in Section~\ref{sec:nogo} uses no empty
frame and remains useful for that reason. Sufficiency uses the genuine
one-sided pseudo laws in Lemma~\ref{lem:covariance}; replacing those by lax
laws could remove the invertible covariance and would require a different
theorem. No extra $\J$-cartesian lift along $\varnothing\to B$ was used.

\section{Uniform receipt framing and its exact scope}\label{sec:receipts}

For a finitary cartesian polynomial monad $T$, a receipt grade over $B$ is
a map $q:Q\to TB$ and has extension $dQ=B+Q$. Native grade maps preserve
$q$. Endpoint transport changes the labels and retains all tags. Cartesian
lift along $v:E\to B$ has tags $Q\times_{TB}TE$. All tags are retained,
including those never emitted by a body.

For any $X$ define
\begin{equation}\label{eq:receipt-frame}
 S_XQ=X\times Q,\qquad q^X(x,r)=T(i_x)q(r),\qquad
 b_{X,Q}:X\times(B+Q)\cong(X\times B)+(X\times Q).
\end{equation}
Use the unchanged context at each $(x,a)$, $b(\id_X\times f)$ on bodies,
and $\id_X\times\beta$ on tag maps. This includes dormant tags and
nullary labels.

\begin{proposition}\label{prop:receipt-positive}
Uniform receipt framing satisfies \textup{(F1)--(F5)} for every frame set.
It is strong for each fixed frame, coherent on the bijection core of
$\Set$, and lax in arbitrary frame functions. Its general parameter
comparison need not be invertible.
\end{proposition}
\begin{proof}
Product/coproduct reassociation gives (F1)--(F2); the tag maps satisfy their
triangle and pentagon by ordinary tuple equality. Endpoint opcartesian
transport only relabels tags, so $S_X$ preserves it. For a context projection
$\pi:E\to B$ there is a grade isomorphism
\begin{equation}\label{eq:receipt-lift}
 X\times(Q\times_{TB}TE)
 \xrightarrow{\sim}(X\times Q)\times_{T(X\times B)}T(X\times E),
 \qquad(x,r,t)\longmapsto((x,r),T(j_x)t),
\end{equation}
where $j_x(e)=(x,e)$. In the inverse direction the polynomial equality
forces the same shape and frame $x$ at every addressed position; remove
these frames. At arity zero the separately retained $(x,r)$ still supplies
$x$. The two functions are inverse and preserve the full receipt labels.
Replacing $\pi$ by any admitted $v$ proves cartesian preservation in (F3).
Naturality is equality of coordinate functions; coefficients at every
address are the original maps.

The receipt unit has no tags, so it is preserved. Accumulation retains the
disjoint union of old and new tags; its original labels are obtained by
transport and $\mu^T$. The distributivity bijection
$(X\times R)+(X\times Q)\cong X\times(R+Q)$ preserves those labels by
naturality of $\mu^T$ and constant-frame insertion. It is the accumulation
isomorphism (F4), on every tag, not just on tags reached by a body.
Its unit and associativity equations send each tag to the same original
stage tag with the same global $x$; the label equations are the monad laws.
(F5) is the definition of $q^X$ and the unary arity of the pure summand.
Theorem~\ref{thm:transfer} proves all mixed comparisons and coherences.
Explicitly a mixed composite with later grade $R$ has tags
$R+(Q\times_{TB}TE)$; after framing its tag comparison is
\[
 (X\times R)+((X\times Q)\times_{T(X\times B)}T(X\times E))
 \cong X\times(R+(Q\times_{TB}TE)).
\]
Every dormant tag is present in this bijection.

For $u:X\to Y$, the vertical parameter factor is $u\times\id_Q$ on tags.
Taking $X=1,Y=2,Q=1$ makes it a map from one tag to two tags, hence not an
isomorphism. Thus (F6) fails and the general comparison is only lax.
\end{proof}

Receipts also verify Proposition~\ref{prop:copower}: total coproducts are
disjoint unions of bases and tags, labelled through the inclusions;
$p,d,U$ preserve them, as does accumulation by its tag formula. For a
vertical map, $d\beta=\id_B+\beta$ is bijective exactly when $\beta$ is.
One may further annotate every receipt by an object of a fixed small
groupoid and allow an annotation arrow on every native tag map. The
extension forgets annotations and is then generally nonfaithful, but is
vertically conservative: a tag bijection with groupoid arrows has an
inverse. All coproduct and transport checks remain valid. This confirms
that the argument used conservativity, not faithfulness of extension.

\subsection{The different failure of projection pullback framing}
There is a tempting but incorrect strong recipe: pull a receipt grade back
along $X\times B\to B$, thereby allowing arbitrary frame labels at its
individual polynomial positions. Even where that projection is admitted,
this can fail strong sequential composition.
Take $T=\List$, $A=B=E=1$, singleton contexts, and $|X|=2$. Let $F$ emit
its unique unary receipt $q$ and let $G$ emit its unique binary receipt $r$.
Their composite has retained tags $r,q$, each finally labelled by the binary
list. Pullback-framing the composite gives $2\cdot2^2=8$ tags. Composing
the separately pullback-framed computations gives four decorations of $r$
and two diagonal decorations inherited from the originally unary $q$, hence
six tags. The missing tags are the two off-diagonal decorations of $q$.
No native grade isomorphism can exist, and singleton contexts cannot repair
that. This is a recipe counterexample; it is distinct from the universal
no-go proved next.

\section{A recipe independent retained receipt obstruction}\label{sec:nogo}

\begin{theorem}[No full Set parameter pseudoaction]\label{thm:nogo}
On the full retained receipt mixed construction for an O1 cartesian
polynomial monad $T$, there is no action satisfying
\textup{(A2)--(A3)} with the stipulated objects $X\times A$, unit comparison
$\st_1\simeq\id$, and pure parameter components $J(u\times\id_A)$.
In particular no choice of framing recipe can provide the full requirement
\textup{(A1)--(A3)} on that construction.
\end{theorem}
\begin{proof}
Let $F:1\to1$ have singleton context and one receipt $q$ with label
$\eta^T_1(*)$; let its body emit that receipt. Put $G=\st_2(F)$ and let
$r:R\to T2$ be its entire retained output receipt grade. There is no
restriction on how the alleged action chose $G$ or its context.
For the two inclusions $i_k:1\to2$, $k=0,1$, invertible parameter
naturality and $\st_1\simeq\id$ give, after the pure unit identifications,
\[
 G\circ J(i_k)\cong J(i_k)\circ F.
\]
Pure precomposition of $G$ restricts its input context but leaves its full
global receipt grade $R$ unchanged: the preceding pure computation has
empty receipt grade. Pure postcomposition of $F$ retains its single tag and
relabels it by $T(i_k)$; crossing the following unit context changes no tag.
Thus the right side has singleton grade labelled $\eta^T_2(k)$.

A mixed isomorphism has invertible context and grade components, because
projection to either component is a functor. Consequently $R\to T2$ must
be isomorphic in $\Set/T2$ both to the singleton labelled $\eta^T_2(0)$
and to the singleton labelled $\eta^T_2(1)$. Such singletons can be
isomorphic only if their labels are equal. They are distinct: cartesianness
of $\eta^T$ at $2\to1$ makes $\eta^T_2$ a pullback of the injective map
$\eta^T_1:1\to T1$, hence injective. This is a contradiction.
\end{proof}
The argument already works for $T=\id$ and uses only two maps between
nonempty finite frames. It does not assume uniform tags, a projection
pullback formula, a fixed shape for the framed contexts, or an empty frame.
It does not exclude strong fixed-frame functors, bijection-core actions,
lax parameter comparisons, or native endpoint squares. In particular those
surviving notions must not be advertised as the forbidden full action.

\begin{corollary}[Invariance of the obstruction]\label{cor:nogo-invariance}
Let $E:H\to H'$ be a biequivalence identity on the pure base objects and
preserving $J$ through coherent invertible comparisons, with a coherent
pseudoinverse. A full action of the prohibited kind on $H'$ would induce
one on $H$. Thus a pure-compatible change to an equivalent presentation
cannot repair Theorem~\ref{thm:nogo}.
\end{corollary}
\begin{proof}
Let $U$ be a pseudoinverse and transport a hypothetical action by
$\st^H_X=U\st^{H'}_XE$. Transport the action constraints using the
unit/counit of the biequivalence and those of the action. The parameter
components are coherently isomorphic to $J_H(u\times\id)$ because $E,U$
preserve pure maps. Conjugating by these component isomorphisms gives the
specified pure components with invertible naturality cells. The axioms
hold by the action axioms and equivalence triangle identities; alternatively
apply the locally faithful equivalence $E$ to each required equality and
cancel the unit/counit comparisons. The resulting action contradicts
Theorem~\ref{thm:nogo}. No frame-preservation structure on $E$ was assumed:
the hypothetical target action itself was transported back.
\end{proof}

\section{Structured equivalences of independent presentations}\label{sec:equivalence}

O1 already proves strong mixed functoriality under its independent parameter
maps, their composition coherence, and local full faithfulness when the
resource and grade functors are fully faithful. Those results are reused
and credited, not counted again as new O2 results. We add essential
surjectivity for a declared class of equivalences, the frame qualification,
and the information-forgetting comparison of the next section.

\begin{definition}[Structured equivalence]\label{def:structured-equivalence}
Over the fixed pure base $\Set$ and a common admitted projection class, take
an O1 parameter map $(F,H,\theta,\delta)$ consisting of a strong
symmetric-rig functor $F:\V\to\V'$, a cartesian polynomial-monad map
$\theta:T\to T'$, a strong contentad functor $H:\Lcat\to\Lcat'$ preserving
specified opcartesian and admitted cartesian transport, and a coherent
extension isomorphism $\delta:d'H\cong d$ satisfying
\[
 \nu'_{HQ}=\theta_{pQ}\nu_Q\delta_Q.
\]
Call it a structured equivalence when $F$ is an equivalence, $\theta$ is
invertible, $H$ is an equivalence over $\Set$ and fibrewise essentially
surjective and fully faithful, and coherent inverse input data have been
chosen so that their units and counits preserve all the indicated rig,
contentad, extension, arity and lift structure.
These are separate-input conditions; no mixed-law comparison is supplied.
\end{definition}
An unstructured equivalence of categories need not preserve augmentation,
occurrence arity, extension bodies or selected marks and is not covered.

\begin{theorem}[Conditional representation invariance]\label{thm:equivalence}
A structured equivalence induces a strong local equivalence
$K:H_{\V,\Lcat,T}\to H_{\V',\Lcat',T'}$, coherent on all endpoint
squares and preserving the designated original-grade and display marks.
Construction before and after this change therefore agrees up to the
specified invertible comparison. These comparisons compose coherently and
preserve the pure inclusion through its unit isomorphisms.
\end{theorem}
\begin{proof}
Apply $F$ coefficientwise, including the global augmentation through the
strong sum and unit isomorphisms, and $H$ to grades. Conjugate the body by
$\delta^{-1}$ and the context-extension identification. This is the O1
mixed functor. Its crossing comparison uses the unique projection-preserving
isomorphism $H(\pi^*Q)\cong\pi^*HQ$, the position bijections of the
cartesian $\theta$, and the strong tensor maps of $F$. The hypotheses match
O1's parameter-map theorem exactly, so strong $j,m$ preservation, all native
squares, the four crossing modifications and their coherences follow from
that proved theorem.

To prove local essential surjectivity, let $(P',Q',f')$ be a target object
$A\to B$. Keep the finite carriers of $P'$ unchanged. Choose for each
coefficient an object $V_{a,p}$ and an isomorphism
$F(V_{a,p})\cong V'_{a,p}$. The strong finite-sum and unit maps identify
the image of the candidate augmentation domain with the target domain.
Fullness of $F$ then lifts the target augmentation, conjugated through those
isomorphisms, to a map $\bigoplus_pV_{a,p}\to I$.
Thus the chosen isomorphisms give an augmented-family isomorphism
$F_*P\cong P'$. Its carrier map is identity, so the admitted projection
condition is unchanged. Choose $Q\in\Lcat_B$ and an isomorphism
$h:HQ\cong Q'$ by fibre essential surjectivity. Conjugating $f'$ by the
context isomorphism, $d'h$ and $\delta$ gives a body $f:e_AP\to dQ$.
The chosen components then give a mixed isomorphism
$K(P,Q,f)\cong(P',Q',f')$ by the very definition of $f$.

For local full faithfulness, a target cell between two image objects has
an unchanged carrier function. Fullness of $F$ lifts each aggregate
coefficient map and faithfulness reflects its augmentation equation.
Fibrewise fullness of $H$ lifts the grade map; faithfulness makes these
lifts unique. Naturality and invertibility of $\delta$ reflect the comma
body equation. The same argument for endpoint squares inserts the canonical
comparisons $F_*(u^*P)\cong u^*F_*P$ and $H(v_!Q)\cong v_!HQ$.
Hence this is a local equivalence on the retained categories, not just a
bijection of closed observations.

Apply O1 to the chosen structured inverse inputs. The coefficientwise and
gradewise units/counits act on the comma bodies by their extension equations,
so they give the corresponding mixed unit/counit isomorphisms. Their
triangle identities follow on the two native components and hence on mixed
cells. Preservation of original-grade projections follows from the defining
triangle of the cartesian comparison; preservation of occurrence marks
follows from the position bijections and the arity equation.

For composable input changes, the two context routes are the strong-rig
coherence for the composite functor. The two grade routes into the same
chosen cartesian lift have the same original projection, so cartesian
uniqueness identifies them. The two body routes are the composite
$\delta$ equation. These equalities prove identity, composition and
triple-composition coherence. Pure computations use singleton coefficient
$I$ and content grade $U$; the input unit comparisons preserve both and
their body. This proves compatibility with $J$ on pure arrows and their
endpoint squares.
\end{proof}

\begin{proposition}[The additional frame condition]\label{prop:equivalence-frame}
Suppose the two inputs have the independently specified framing data of
Definition~\ref{def:frame-data}. Give coherent one-sided isomorphisms
$\chi_{X,Q}:HS_XQ\cong S'_XHQ$, compatible with unit, accumulation,
transport and all frame parameters, and satisfying
\begin{equation}\label{eq:frame-morphism}
 d'\chi_{X,Q}\,\delta^{-1}_{S_XQ}b_{X,Q}
   =b'_{X,HQ}(\id_X\times\delta_Q^{-1}).
\end{equation}
Then the induced mixed comparison intertwines the constructed frame
actions, including their sequential and parameter constraints. This extra
condition cannot be inferred from an equivalence of underlying categories.
\end{proposition}
\begin{proof}
Both framed context constructions retain the same carriers; their
coefficient comparison is the strong structure map of $F$, independent of
$x$. On grades use $\chi$. Equation~\eqref{eq:frame-morphism} is precisely
the mixed body equation. For the crossing, both grade paths preserve the
same cartesian projection, so uniqueness equates them. The polynomial
position maps commute with constant-frame insertion because
\[
 \theta_{X\times B}T(i_x)=T'(i_x)\theta_B.
\]
Thus their occurrence coefficients are the same $F$-images on the same
addressed positions, identified by strong-rig coherence. Compatibility
with action and sequential constraints follows respectively from the
specified one-sided $\chi$ coherences and the unchanged coefficient
coherences. This is a derived mixed compatibility, not an extra mixed-law
assumption. In the constant-polynomial or fixed-polynomial canonical cases,
the body part is obtained directly by this naturality equation whenever
the stated one-sided grade comparison is available.
\end{proof}

The existence/nonexistence statement in Corollary~\ref{cor:nogo-invariance}
is different: it transports a hypothetical action along a pure-compatible
biequivalence, and needs no previously chosen $\chi$. The present proposition
compares two already chosen actions and so specifies exactly how they are
to match.

\section{Arity shadows and probability law quotients}\label{sec:probability}

The representation comparison can exercise the full chosen constant
polynomial grade family, not just isomorphic relabellings or its unary
subfragment. We first isolate the content information forgotten, then the
sample information forgotten. The ordinary finite-kernel case will be a
control of the same formulas.

\subsection{A strong arity shadow for every O1 content input}
Let $H_{\V,T}$ denote O1 with the same contexts and terminal-grade content
$d_B(*)=TB$. Define
\begin{equation}\label{eq:arity-shadow}
 E_\nu:H_{\V,\Lcat,T}\longrightarrow H_{\V,T},\qquad
 (P,Q,f)\longmapsto(P,\nu_Qf),\quad(\alpha,\beta)\longmapsto\alpha.
\end{equation}
This is not generally an O1 parameter map of Definition
\ref{def:structured-equivalence}, since $\nu$ need not be invertible.

\begin{theorem}[Strong arity shadow]\label{thm:arity-shadow}
The map~\eqref{eq:arity-shadow} is strong for units and sequential
composition, preserving all native context cells, all pure endpoint squares,
and the pure inclusion. Wherever the source has framing satisfying
\eqref{eq:frame-arity}, it preserves fixed-frame comparisons and their
parameter squares. A lax source parameter comparison may become invertible
in the target; this does not make the source action pseudonatural.
\end{theorem}
\begin{proof}
For a native cell, $d\beta f=f'e\alpha$ and naturality of $\nu$ give
$\nu_Qf=\nu_{Q'}f'e\alpha$. For an endpoint square,
\[
 \nu_{v_!Q}d\kappa_{v,Q}=Tv\,\nu_Q
\]
gives the fixed-$T$ square with the same context map $\alpha$.
Thus cells, their composition, and pure transports are preserved.

Let $F=(P,Q,f):A\to B$ and $G=(R,S,g):B\to E$. For an initial witness
$p$, set $t=\nu_Qf(p)\in TB$. The original mixed composite context is
one independent $R$-witness for each addressed occurrence of $t$, together
with $p$ and its coefficient. This is exactly the fixed-$T$ composite
context, including augmentations and actual native coefficient maps.
Write $\widetilde t\in T(e_BR)$ for the resulting occurrence-decorated
point. The fixed-$T$ composite body is
\begin{equation}\label{eq:shadow-composite}
 \mu^T_E\,T(\nu_Sg)(\widetilde t).
\end{equation}
In the source, the display equation for the crossing identifies the arity
of the lifted body with $\widetilde t$. Naturality for transport along
$g$ identifies its subsequent arity with the corresponding transported
$T$-point. Finally the content arity accumulation equation
$\nu\mu^L=\mu^T T\nu\,\nu$ gives exactly
\eqref{eq:shadow-composite}. No inverse to $\nu$ is used.
The multiplication comparison is therefore the canonical identity on the
common addressed context, with equal target bodies. It is invertible.

The arity-unit equation $\nu_U\eta^L=\eta^T$ supplies the pure and mixed
unit comparisons. For triple composition, both context routes are the
same O1 addressed dependent-tensor map on the same primitive coefficients;
these actual maps are retained by $E_\nu$. Applying the arity accumulation
equation twice makes the two body routes the two sides of the monad
associativity equation. Unit coherence is the corresponding monad unit
equation and unchanged context unitor. Native naturality holds because the
context aggregate maps use exactly those same addressed tuples. This
proves strong coherence, not only equality of final laws.

Lastly~\eqref{eq:frame-arity} says exactly that applying $E_\nu$ before
or after framing gives the same fixed-$T$ body. All context components
are unchanged. The action and sequential comparisons are consequently the
same addressed context maps. For a source parameter square, arity
naturality gives the target parameter square, while any noninvertible
content-grade factor is forgotten. This proves the qualification.
\end{proof}

\subsection{The local reflection for the full constant collection family}
Let $\mathbb L$ be the terminal collection, with one shape in each arity.
Then $F_{\mathbb L}=\List$, and the unique grade map
$!_C:C\to\mathbb L$ has extension $\nu_C$.
For a fixed-List history $(P,h:e_AP\to\List B)$ set
\[
 R(P,h)=(P,\mathbb L,h),\qquad R(\alpha)=(\alpha,\id_{\mathbb L}).
\]
This is a section of $E_\nu$ on each hom category.

\begin{proposition}[Local reflection and a lax section]\label{prop:arity-reflection}
For every $A,B$ and every admitted rig, $E_\nu\dashv R$ locally and
$E_\nu R=\id$. These local comparisons respect pure endpoints and the full
Set frame action. The sequential section $R$ is only lax; both its unit
and its multiplication comparisons can be noninvertible.
\end{proposition}
\begin{proof}
An arrow $(P,C,f)\to R(P',h)$ consists of a context arrow $\alpha$ and a
grade map $C\to\mathbb L$. The latter is unique, and the body equation is
exactly $\nu_Cf=h e\alpha$. This gives the natural bijection
\[
 \Hom_H((P,C,f),R(P',h))
 \cong\Hom_{H_{\V,\List}}((P,\nu_Cf),(P',h)).
\]
The local unit is $(\id_P,!_C)$ and the counit is identity.
Endpoint transports are constant on collection grades and $\nu$ is natural
in the endpoint, proving the endpoint assertion.

The sequential comparison $R(G)\circ R(F)\to R(G\circ F)$ retains the
same context and uses the grade map
$\Sub(\mathbb L,\mathbb L)\to\mathbb L$ that forgets the ordered
partition of the final leaves and remembers only their total arity.
Its extension body is list flattening. This map is generally not an
isomorphism: many distinct ordered inner-length lists have the same sum.
The unit comparison uses $\mathbb I\to\mathbb L$, selecting the unary
shape, and is also not an isomorphism. Both routes through a nested triple
forget its tree partition and retain the same ordered leaf list, so their
grade maps are equal. Their context maps are the same inherited O1
associator, and their bodies are the equal list flattenings. Unit diagrams
follow from singleton-list insertion and deletion. This proves lax
sequential coherence with its exact direction.

Framing commutes with $E_\nu$ by forgetting the shape after tagging leaves,
and with $R$ because $\mathbb L$ is unchanged. These identifications
preserve the parameter, unit, action-associator and sequential comparisons:
all their bodies tag the same ordered leaves and their grade maps are the
ones just described. Thus a locally fully faithful right adjoint need not
be a strong sequential section.
\end{proof}

\subsection{Raw finite probability histories and independent list laws}
Now choose the discrete nonnegative-real semiring rig, whose objects are
nonnegative numbers, with $\oplus=+$, $\otimes=\cdot$, $I=1$, and only
identity arrows. Choose $\J$ to contain every finite witness projection;
for definiteness take all set maps. The normalized contexts then consist
of finite sets $P_a$ with $p_a(x)\ge0$ and $\sum_xp_a(x)=1$.
Zero-weight witnesses remain native data. A native context arrow is a
function $r_a:P_a\to P'_a$ such that
\[
 p'_a(y)=\sum_{r_a(x)=y}p_a(x).
\]
Thus all finite probability histories used below are admitted, not merely
a subset accidentally allowed by an unspecified $\J$.

In the terminal-grade List model, a history also has an output list
$h_a:P_a\to\List B$. Native arrows obey $h_a=h'_ar_a$.
For an endpoint square over $u:A\to A'$ and $v:B\to B'$, the output
equation is $\List(v)h_a=h'_{u(a)}r_a$, with the same weight pushforward.
The composition of $(P,p,h):A\to B$ and $(Q,q,k):B\to E$ has witnesses
\begin{equation}\label{eq:list-histories}
 (x,(y_i)_{i=1}^n),\quad h_a(x)=[b_1,\ldots,b_n],\quad
 y_i\in Q_{b_i},
\end{equation}
weight $p_a(x)\prod_iq_{b_i}(y_i)$, and output
$k_{b_1}(y_1)\mathbin{+\!+}\cdots\mathbin{+\!+}k_{b_n}(y_n)$.
These are exactly the O1 occurrence witnesses and coefficients, not a
separately imposed sampling operation. Repeated occurrences of the same
$b$ use independently chosen witnesses; sharing a sample would be a
different semantics.

Let $\mathsf K_{\List}(A,B)$ be the discrete category of maps
$A\to\mathcal D_{\mathrm{fin}}(\List B)$, where $\mathcal D_{\mathrm{fin}}$
means normalized finitely supported laws. Define composition by
\begin{equation}\label{eq:list-kernel}
 (\ell\star k)_a(w)=
 \sum_{\mathbf b=[b_1,\ldots,b_n]}k_a(\mathbf b)
 \sum_{w_1+\!+\cdots+\!+w_n=w}\prod_{i=1}^n\ell_{b_i}(w_i).
\end{equation}
Only supported lists and supported continuation lists contribute; thus both
sums are finite. At $n=0$ the empty product is $1$ and the output is the
empty list. The unit is the Dirac law at the singleton list.
Formula~\eqref{eq:list-kernel} defines a normalized finite-support law:
for each supported initial list the sum over all continuation choices is
$\prod_i\sum_{w_i}\ell_{b_i}(w_i)=1$, and the outer masses sum to $1$.
For triple composition both bracketings enumerate the same finite tree of
an initial list, continuation lists at its addresses, and final lists at
their addresses. The weight of a tree is the product of all its factors;
both sums group the same trees by the same flattened final list. Finite
distributivity and associative concatenation prove associativity. Singleton
Dirac lists prove the unit laws. This is the independent-sampling Kleisli
composition for $\mathcal D_{\mathrm{fin}}\List$, not ordinary one-output
kernel substitution.

Define the list-law map
\begin{equation}\label{eq:list-law}
 L(P,p,h)_a(w)=\sum_{x:\,h_a(x)=w}p_a(x),\qquad
 G=L E_\nu.
\end{equation}
Here $G$ is defined on the full constant-polynomial-grade probability
fragment, not just on grade $\mathbb L$.

\begin{theorem}[Strong list law quotient]\label{thm:law}
The maps $L$ and $G$ are strong for sequential composition and units, on
all source native cells and endpoint squares. They preserve the full
deterministic Set frame action. The map $G$ is locally surjective on objects
but is neither locally full nor locally faithful, and does not reflect
invertibility. It forgets actual shape and sample-witness information.
\end{theorem}
\begin{proof}
A native history arrow preserves~\eqref{eq:list-law} since
\[
 \sum_{x:h_a(x)=w}p_a(x)
 =\sum_{y:h'_a(y)=w}\sum_{x:r_a(x)=y}p_a(x)
 =\sum_{y:h'_a(y)=w}p'_a(y).
\]
It therefore maps to the unique identity at that law. For an endpoint
square the same finite sum gives
$(\List v)_*L(P)_a=L(P')_{u(a)}$, the exact deterministic pushforward square.
This equation is preserved under composed endpoint maps.

Expanding the composite history~\eqref{eq:list-histories} gives
\[
 \sum_{x}p_a(x)
 \sum_{(y_i):\, k_{b_1}(y_1)+\!+\cdots+\!+k_{b_n}(y_n)=w}
                  \prod_iq_{b_i}(y_i).
\]
Group the outer summands by $h_a(x)=\mathbf b$ and inner summands by
$k_{b_i}(y_i)=w_i$. Finite distributivity changes each inner group to
$\prod_i\sum_{y_i:k_{b_i}(y_i)=w_i}q_{b_i}(y_i)$.
The result is exactly~\eqref{eq:list-kernel} for $L(P)$ and $L(Q)$.
This includes the empty list and zero weights. The unit has singleton
Dirac law. Hence comparison maps in the discrete target are equalities.
Their coherence follows from the proved object equalities, the common
addressed-tree associator and the finite-tree associativity calculation
above; only after proving those equalities is uniqueness of target cells
applicable. Theorem~\ref{thm:arity-shadow} proves the assertions for $G$.

For fixed frame $z\in Z$, let $\lambda_z:\List B\to\List(Z\times B)$
map $[b_i]_i$ to $[(z,b_i)]_i$. Frame a law by
\[
 (\st_Zk)_{(z,a)}=(\lambda_z)_*k_a.
\]
The framed history has the same weights and output $\lambda_zh_a$, so
$L\st_Z=\st_ZL$ exactly. The equality for $G$ also uses
$\nu b=\operatorname{st}^{\List}(\id\times\nu)$.
Concatenation commutes with $\lambda_z$, hence these equalities preserve
strong sequential comparisons. For $t:Z\to Z'$,
$\List(t\times\id)\lambda_z=\lambda_{t(z)}$; this proves arbitrary
parameter naturality. Endpoint naturality, action units and associators
are the same identities on list labels. Thus framing is fully preserved,
not merely checked on a marginal of a single output.

For surjectivity, realize a law $k_a$ by its finite positive support,
weight $k_a(w)$ at $w$, output $w$, and grade $\mathbb L$.
For nonfaithfulness, take two witnesses of weight $1/2$, both with the same
unary output list $[b]$. Identity and transposition are distinct native
endomorphisms and have the same law image. For nonfullness, a one-witness
probability-one history and the preceding two-half history have the same
law. There is no native arrow from the former to the latter: a deterministic
map pushes mass one to one point, never to two halves. Yet the discrete
target has its identity arrow. The reverse two-to-one collapse is native
and noninvertible but maps to an identity, disproving reflection of
invertibility. All examples can use the same grade $\mathbb I$ or
$\mathbb L$, so no grade issue hides the witnesses.
Finally $!_C:C\to\mathbb L$ forgets shape identity, including dormant
shapes. For example two distinct unary shapes can be merged to the one
unary shape without changing the ordered output-list law. This is an
additional loss beyond sample aggregation.
\end{proof}

\subsection{The exact raw quotient and the positivity boundary}
For a category $D$, write $\pi_0D$ for the discrete category of connected
components of its underlying undirected graph. This does not invert native
arrows to a groupoid or assert contractibility of any component.

\begin{theorem}[Local quotient and positive support adjunction]\label{thm:quotient}
For all $A,B$, the map $G$ induces a bijection
\[
 \pi_0 H_{\mathrm{prob},\C}(A,B)
       \cong\mathsf K_{\List}(A,B).
\]
Consequently every functor from that hom category to a discrete category
factors uniquely through $G$. The quotient respects sequential composition,
all pure endpoints, and the full frame action.

On the subcategory of positive-weight witnesses, $G$ has a local right
adjoint $S$ given by the support sampler at grade $\mathbb L$, with $GS=\id$.
This $S$ is globally only lax for sequential composition, and its unit
comparison in the full constant-grade model is noninvertible. The same
support sampler is not a local right adjoint on raw histories with zero
witnesses.
\end{theorem}
\begin{proof}
Every native arrow preserves $G$, so connected objects have the same law.
For $F=(P,C,f)$, the native unit $(\id_P,!_C)$ connects it to
$R E_\nu F=(P,\mathbb L,\nu_Cf)$.
Delete all zero-weight witnesses of this latter history, obtaining $P^+$.
Its inclusion $P^+\to P$ is native: omitted target points have weight
zero, and outputs are unchanged. Every remaining witness has positive
mass, so its output lies in the positive law support. Mapping it to that
output gives a native arrow to $S(GF)$, with the correct aggregated weight
by~\eqref{eq:list-law}. Thus there is a zigzag
\[
 F\longrightarrow R E_\nu F
       \longleftarrow (R E_\nu F)^+
       \longrightarrow S(GF).
\]
Two objects of the same law are connected through this common support
sampler. Surjectivity was proved in Theorem~\ref{thm:law}, giving the
bijection. A functor to a discrete category is constant on each zigzag;
this proves the unique factorization, on objects and arrows. Sequential
composition, endpoint transport and framing are functorial on native cells,
so they carry zigzags to zigzags. Their quotient operations are identified
with the stated target operations by Theorem~\ref{thm:law}.

Now assume every source witness has positive weight. A native arrow
$F\to S(k)$ must have the unique shape map $C\to\mathbb L$ and must send
a witness $x$ to the list $\nu_Cf(x)$, since the sampler outputs its carrier
list. This list lies in the support exactly as needed when $G(F)=k$;
its pushforward weights equal $k$ by definition. There is therefore exactly
one such arrow when $G(F)=k$ and none otherwise. These are precisely the
hom sets in the discrete law category, giving the natural adjunction
$G\dashv S$ locally. Its counit is identity and its unit merges equal-list
witnesses and forgets their collection shapes.

For sequential composition the sampler comparison
$S(\ell)\circ S(k)\to S(\ell\star k)$ maps a complete tuple of sampled
lists to its concatenated final list. Positive products ensure this final
list is in the positive support. The pushforward weight equation is
\eqref{eq:list-kernel}. Its grade component is
$\Sub(\mathbb L,\mathbb L)\to\mathbb L$. Both associative routes send
an addressed tree of samples to the same final concatenation; their weight
equation is the same finite sum/product expansion and their grade map
forgets the same tree partition. The unit comparison has grade
$\mathbb I\to\mathbb L$. This proves lax associativity and unit coherence.
It is not strong: already a fair first choice followed by a constant unary
output merges two positive samples, and the full-model grade unit map is
not invertible. For framing, deleting the constant frame from each list
is inverse to $\lambda_z$ on its image. It therefore gives the canonical
bijection between the support of $k_a$ with frame $z$ and the support of
$(\st_Zk)_{(z,a)}$. This proves the sampler's frame compatibility, including
nullary lists, by an actual support bijection.

For the raw-zero counterexample, use one positive witness of output $[b_0]$
and weight $1$, and one zero witness of output $[b_1]$, with $b_1\ne b_0$.
The law support is $\{[b_0]\}$, and no output-preserving map can send the
second witness there. Thus even the proposed unit $F\to S(GF)$ fails.
The raw connected-component quotient is unaffected because it used the
reverse inclusion from the positive subhistory, not that nonexistent map.
\end{proof}

\subsection{The unary kernel control and the receipt comparison}
Restrict constant collection content to $\mathbb I$ and bodies with exactly
one output leaf. This grade is closed under substitution, and its extension
is $B$. It is the fixed-$T=\id$ history model inside the chosen positive
fragment. The list-law formula reduces to ordinary finite kernels
$k:A\to\mathcal D_{\mathrm{fin}}(B)$ with
\[
 (\ell k)_a(c)=\sum_b k_a(b)\ell_b(c),\qquad
 E(P)_a(b)=\sum_{x:h_a(x)=b}p_a(x).
\]
Its support-sampler compositor sends $(b,c)\mapsto c$ and is generally
noninvertible; its pure unit comparison is invertible in this unary model.
On positive witnesses the local adjunction $E\dashv S$ holds, while the
same raw-zero counterexample blocks it on all histories. Its nonfaithfulness,
nonfullness and exact connected-component quotient are the unary
specializations above. This is a useful control, not a replacement for the
full-grade list-law result.

For retained receipts, Theorem~\ref{thm:arity-shadow} instead gives the
strong map $(P,Q,f)\mapsto(P,\nu_Qf)$ to fixed-$T$ histories. In the
probability $T=\id$ fragment one can then take the ordinary kernel law;
for $T=\List$ one can take the list law just constructed. The first map is
not generally full or faithful. For example, two different maps between
dormant tags can have the same context component, while an ordinary return
and a used unary receipt can have the same arity body without a grade map
from the receipt grade to the empty return grade. Thus no receipt-level
support adjunction is asserted, even with positive sample weights.
The composite forgets global receipt provenance as well as sample history.
It preserves the strong fixed-frame comparisons and maps lax source
parameter comparisons to target invertible comparisons. This is consistent
with Theorem~\ref{thm:nogo}: neither step is a pure-compatible biequivalence
of the retained receipt presentation.

\subsection{Observations preserved and observations excluded}
For any bounded function $\varphi:\List B\to\mathbb R$, the exact identity
\[
 \sum_xp_a(x)\varphi(\nu_Cf_a(x))
       =\sum_wG(F)_a(w)\varphi(w)
\]
proves preservation of all final ordered-list events and bounded tests.
The unary case gives the usual final-output expectation identity.
Tests of a private witness, a collection shape, a dormant receipt, or an
annotation are not covered and in general are not preserved.

Nor does equality of these laws permit adding hidden shared-sample
operations. A box retaining one latent fair bit has the same one-call law
as fresh fair sampling, but equality of two calls has probability $1$ in
the first mechanism and $1/2$ in the second. The sequential operation in
\eqref{eq:list-histories} uses independent occurrences; the shared-state
operation is absent. The example is therefore a boundary on extending the
observation language, not a counterexample to the strong quotient theorem.
No higher-order or unrestricted contextual equivalence follows.

\section{Preservation and loss ledger}\label{sec:ledger}

\begin{longtable}{@{}p{.18\linewidth}p{.35\linewidth}p{.39\linewidth}@{}}
\toprule
Feature&Structured input equivalence&Full constant polynomial probability law map $G$\\
\midrule\endhead
Interfaces&Same pure base; no hidden enlargement.&Same $A,B$; laws are on
ordered lists of $B$, explicitly exposing the output regime.\\
Mixed objects&Locally essentially surjective and fully faithful.&Surjective
onto finite laws; distinct shapes and witnesses can have the same law.\\
Native cells&Bijections on hom sets; coefficient and grade maps retained
up to the specified equivalence.&All cells map to equality. Parallel
permutations collapse; some target equalities have no direct source arrow.\\
Endpoint squares&Transport-compatible equivalence on actual squares.&Every
source square gives exact $\List(v)$ pushforward. Not every target square
lifts.\\
Pure maps&Preserved through coherent rig and content unit comparisons.&Sent
to Dirac singleton lists, exactly the target pure inclusion.\\
Observations&Transport of the declared marks and observation structure;
not arbitrary unmarked observations.&Final ordered-output laws and bounded
tests preserved; shape, sample, dormant-grade and latent-state tests lost.\\
Sequential composition&Strong comparison with coherent strong inverse.&Strong
equality for independent occurrence sampling and concatenation. Not a
shared-sampling theorem.\\
External action&Chosen actions match under the one-sided condition of
Proposition~\ref{prop:equivalence-frame}.&Full Set action exactly preserved,
including arbitrary frame parameters and nonunit frames.\\
Reflection&Invertibility and native equality reflected locally.&Neither
native equality nor invertibility reflected; a collapse can become identity.\\
Quotient or inverse&Structured quasi-inverse, not information erasure.&Raw
local quotient is $\pi_0$. A support right adjoint exists only on positive
witnesses and is sequentially lax.\\
\bottomrule
\end{longtable}

The intermediate arity shadow $E_\nu$ retains every context coefficient map
and witness but can forget content-grade distinctions. For constant
collections it is a local reflection to List histories even for nonnumerical
rigs; its right adjoint is not a strong sequential section. The law map $L$
then forgets probability-witness distinctions. This factorization identifies
the two losses rather than attributing them both to an equivalence.

For receipts the additional framing ledger is essential. Uniform framing
preserves fixed-$X$ sequential structure and all dormant tags, and has
invertible parameter comparisons for bijections. Arbitrary parameter maps
have only lax grade maps $u\times\id_Q$. The stronger no-go excludes any
replacement with full invertible parameter comparisons and the stipulated
pure components. Pure-compatible biequivalence preserves that obstruction;
arity/law quotienting need not. In particular a fully framed target after
forgetting provenance is not evidence of full framing before the quotient.

\section{Assumptions sources and acceptance crosswalk}\label{sec:acceptance}

\subsection{Versioned assumption and result ledger}
The source baseline is the O1 \emph{General Independent Interaction
Criterion} dated 1 November 2026~\cite{O1}, not just its earlier indexed
receipt predecessor~\cite{Receipt}. The ambient relative module,
augmented-rig context doctrine, display-based crossing, its four
modifications and ten coherence diagrams, and strong O1 parameter theorem
are inherited with their original hypotheses. They are not claimed newly
proved here. The present report supplies the exact empty-fibre O2 action
criterion, the explicit positive action checks, the one-sided transfer
mechanism, the two distinct receipt boundaries,
the declared equivalence extension, and the exact information-loss
comparisons at the stated scope. No priority or novelty certification is
claimed for these elementary or structural arguments.

The independent assumptions are separated as follows.
\begin{enumerate}
\item The exact criterion applies to the entire canonical O1 mixed
construction with its genuine one-sided pseudo structure. It includes empty
frames/interfaces and requires the specified pure parameter components.
Equivalence of all $(!_B)_!$ is necessary and sufficient within this class,
not automatically true and not asserted sufficient for arbitrary
bicategories.
\item The explicit constant-polynomial instance uses any admitted O1 rig and
$\J$, with $T=\List$. All native shape maps preserve arity.
Arbitrary deletion of a position,
correlation of distinct occurrences, or alteration of the original O1
projection marks is not silently added.
\item For an already specified framing, the transfer theorem needs
(F1)--(F5), and full parameter pseudonaturality needs (F6). Under the exact
empty-fibre criterion, Section~\ref{sec:exact} constructs all of these
after a structured change of presentation. O1 alone, without that criterion,
does not imply existence of a full action. The finite-copower criterion
still proves only the fixed-frame/bijection version without (F6).
\item The receipt no-go uses global grade retention, the given pure
inclusion and pure parameter components, the unit comparison, and
invertible horizontal parameter comparisons. Changing one of these is a
change of requirement, not a counterexample to its proof.
\item Representation equivalence means the structured class of Definition
\ref{def:structured-equivalence}. Preserving a separately chosen action
requires Proposition~\ref{prop:equivalence-frame}; mere equality of laws
never supplies such an equivalence.
\item The probability theorem specializes to normalized nonnegative
coefficients and a projection class admitting all its finite witnesses
(here all maps). It retains zeros for the quotient theorem, restricts to
positive weights only for the sampler adjunction, and uses independent
samples at addressed occurrences. Its laws are on finite ordered output
lists; ordinary kernels are its unary control.
\end{enumerate}

\subsection{Relation to existing structural baselines}
The parameter requirement is standard for a pseudofunctorial bicategorical
action. Paquet and Saville~\cite{PaquetSaville}, Definition 10 in the cited
author version, take an action to include a pseudofunctor
$\mathcal B\times\mathcal H\to\mathcal H$ and coherent action constraints;
their Theorem 1 relates suitable Kleisli actions and strengths.
For such a pseudofunctor, the two decompositions of $(u,F)$ give invertible
compositor comparisons between
$\mathcal A(u,\id_B)\circ\mathcal A(\id_X,F)$ and
$\mathcal A(\id_Y,F)\circ\mathcal A(u,\id_A)$.
When the parameter components are the stipulated pure maps, this is
\eqref{eq:parameter}. Thus its invertibility is not an extra condition
invented to rule out receipts. Invertible comparison cells do not require
the body map $b_{X,Q}$ itself to be bijective. We do not assert that every
O1 retained construction has already been identified with the paper's
Kleisli setup, nor use its correspondence theorem without that translation.

Shulman's monoidal-fibration framework is an important existing
baseline~\cite{Shulman}. Definition 12.1 includes monoidal structure and
preservation of cartesian arrows, not just an underlying fibration.
Theorem 12.7 relates monoidal fibrations over cartesian bases to indexed
monoidal categories; Theorems 14.2 and 14.9 construct framed bicategories
and functorial comparisons from suitable monoidal bifibrations. These
results do not imply that an arbitrary O1 contentad already has the
additional framing structure used here. A direct reduction would require
matching the retained relative modules and verifying those monoidal and
base-change hypotheses. We use the explicit proofs above, and make no
claim of independence from or priority over the established framework.
The distinction between horizontal pseudonaturality and a native square
is material to the receipt obstruction.

\subsection{O2 acceptance contribution and open programme}
\begin{longtable}{@{}p{.29\linewidth}p{.65\linewidth}@{}}
\toprule
Contract requirement&Evidence and exact status\\
\midrule\endhead
Declared external monoidal structure&$(\Set,\times,1)$ in Section
\ref{sec:scope}, explicitly distinct from relative sequencing and internal
resource tensor. No binary computation tensor is asserted.\\
Exact existence test on the O1 family&Theorem~\ref{thm:main} and Section
\ref{sec:exact}: full pure-compatible Set framing exists exactly when
every $(!_B)_!$ is an equivalence. All-map covariance is derived from
one-sided pseudo validity, including its body equation; no
$!_B\in\J$ hypothesis is imposed.\\
Chosen nontrivial action&Theorem~\ref{thm:main} and the complete component
proof in Section~\ref{sec:constant}; constant ordered polynomial content
satisfies the criterion with all admitted rig contexts and native cells.\\
Pure maps and native cells&Endpoint equations, aggregate coefficient maps
and arbitrary frame-parameter invertible comparisons are proved in Section
\ref{sec:constant}. This meets the strongest parameter reading, not merely
the weaker receipt result.\\
Action unit and associativity&Product label comparisons and their triangle,
pentagon, parameter and sequential compatibility in Section
\ref{sec:constant}.\\
Compatibility with $j,m$ and mixed interaction&Explicit sequential comparison
\eqref{eq:strong-seq}, the restricted crossing comparison and all four
modification checks; Theorem~\ref{thm:transfer} derives the general mixed
comparison from separate one-sided data.\\
Nonunit effectful test&Example~\ref{ex:effect} uses two frame states and a
binary effect shape with actual dependent occurrence witnesses and resources.\\
Conditional representation invariance&Theorem~\ref{thm:equivalence}, including
local essential surjectivity, structured inverse, composed comparisons,
marks and pure compatibility; Proposition~\ref{prop:equivalence-frame}
states the separate frame condition. The intrinsic empty-hom criterion is
also invariant under pure-compatible biequivalence by Section
\ref{sec:exact}.\\
Genuinely different representation&Theorems~\ref{thm:arity-shadow},
\ref{thm:law} and \ref{thm:quotient}: full collection grades and retained
probability histories to extensional laws on ordered lists, with a local
reflection followed by a nonfull/nonfaithful quotient.\\
Preservation and loss ledger&Section~\ref{sec:ledger}; raw zeros, dormant
shapes/receipts, noninvertible cells, positive-only support adjunction, and
independent-sampling observation boundary are explicit.\\
Boundaries and no overclaim&The six-versus-eight recipe failure differs
from Theorem~\ref{thm:nogo}'s universal two-inclusion obstruction; neither
silently replaces the positive full-Set theorem.\\
\bottomrule
\end{longtable}

The result is offered as an O2 acceptance candidate with an exact action
criterion on the entire canonical O1 family and a fully worked nontrivial
positive family. Formal acceptance should record the theorem's quantifiers,
especially empty frames/interfaces, pure parameter components and genuine
one-sided pseudo validity, as well as the probability comparison's observation
regime. The report does not label receipt framing as full Set pseudonaturality
or claim that every O1 input satisfies the criterion. The positive branch
meets the strongest pure-parameter reading; a weaker lax-only requirement
is not substituted for the mandatory result.

The whole expanded research contract remains \emph{incomplete}. This report
does not discharge the O3 plural-language and QBS comparison, the O4 language
metatheory, O5 higher-order/contextual soundness or adequacy, the full O6
application sheets, or the integrated O7--O8 obligations. It does not establish
closed structure, arbitrary state encapsulation, feedback, scheduling,
physical resource copying, partial recovery, or unrestricted presentation
independence. Fuller motivating applications remain later work, rather than
being inferred from structural fixtures. These boundaries do not weaken the
O2 theorem just proved, but they prevent acceptance of that theorem from
being mistaken for completion of the entire programme.

\appendix
\section{Reproducible checks and source only verification}\label{app:checks}
The companion bundle contains this final TeX, the exact standard-library
check scripts, their output, a README, a manifest and SHA-256 hashes. It
excludes primary-source PDFs, private working notes and raw scratch files.
No TeX engine was run, as authorized; no compiled PDF or visual-layout
verification is claimed.

\begin{itemize}
\item \texttt{check\_framing\_review.py} verifies 540 finite uniform-decoration
fixtures with frame sizes zero through three, arities zero through three,
and two labels whose witness-fibre sizes range from zero through two. It
also checks the six-versus-eight pullback recipe count, the one-versus-two
uniform-parameter count, and the two nonisomorphic singleton labels used by
the recipe-independent no-go. These are tests of explicit element formulas,
not proofs of universal coherence or of the universal obstruction.
\item \texttt{check\_representation.py} verifies 4,650 exact rational unary
composition cases, the nonfaithfulness/nonfullness examples, the raw-zero
support obstruction and nonunit frames. It additionally checks 216 List-law
composition cases with repeated occurrences, nullary lists and zero witnesses.
The list suite distinguishes independent list substitution from ordinary
kernel matrix multiplication.
\item \texttt{check\_tex\_source.py} checks balanced TeX groups and environments,
unique labels, resolved local references/citations, and source-integrity
conditions without compiling. This is syntactic evidence only and does not
certify mathematical truth or typeset layout.
\end{itemize}
All unbounded claims rest on the proofs in the report and the explicitly
identified O1 dependency, not on finite test counts. Independent mathematical
review checks the distinctions among full parameter pseudonaturality,
fixed-frame strength and lax comparisons, and among equivalence, reflection
and quotient. Source checking is not proof-assistant certification.

\begin{thebibliography}{9}
\bibitem{Contract}
\emph{Expanded Research Contract v0 1}. Project contract, 1 November 2026,
O2 and its acceptance requirements. \bibitem{O1}
\emph{General Independent Interaction Criterion}. Project proof dossier,
1 November 2026. General augmented-rig and relative-occurrence-display
criterion; strong input changes, retained endpoint cells and examples.
Source \texttt{General\_Independent\_Interaction\_Criterion.tex}.
\bibitem{Receipt}
\emph{Indexed Receipt Interaction Theorem}. Project proof dossier,
1 November 2026. Scoped predecessor, not substituted for the general O1
criterion. \bibitem{Shulman}
M. Shulman. \emph{Framed bicategories and monoidal fibrations}.
Definition 12.1 and Theorems 12.7, 14.2 and 14.9.
\url{https://math.uchicago.edu/~may/PEOPLE/MIKE/framed.pdf}.
Primary source checked 1 November 2026.

\bibitem{PaquetSaville}
H. Paquet and P. Saville.
\emph{Strong pseudomonads and premonoidal bicategories}.
Author-hosted 15-page version, Definitions 8--10 and Theorem 1.
\url{https://www.cs.ox.ac.uk/people/hugo.paquet/strong.pdf}.
Primary source checked 1 November 2026; numbering is version-specific.
\end{thebibliography}
\end{document}
